\chapter{相变、界面与外场平衡}

\section{选择状态变量与守恒的能量}
\label{th:coordinates}
本讲将前面的状态函数方法用于界面和外场。简单可压缩物质的体积功为$-p\dd V$；界面变化产生表面功，外场还会引入势能，须先写清系统所含的能量与全部广义功。第~\ref{th:energyidentity}节已用链式法则证明$U_V|_T=-p+Tp_T|_V$，其中等温与等熵偏导取沿不同的状态方向；改用新的坐标时，仍须分别固定其余变量。

体积偏导只确定$U$的体积依赖，积分后留下的函数$\Phi(T)$还要由定容热容确定。下面先处理总量与摩尔量，再把同一方法用于表面面积及外场参数。

\subsection{广义坐标与相应的功}
\label{th:generalizedwork}
体积、界面面积和外场参数改变时，能量中出现的功项各有不同。把一个可控参数记为 $a$，其共轭广义力记为 $Y$，对固定粒子数的平衡态可写
\[
 \dd U=T\dd S-p\dd V+Y\dd a.
\]
这里 $Y\dd a$ 按外界对系统做功为正。界面以 $a=\Sigma$ 为坐标时 $Y=\sigma$；改变均匀重力场强度时，若系统能量计入重力势能，$a=g$，则 $Y$ 是全部粒子的质量乘高度之和的平均值。共轭量的符号由所选参数及功的约定确定。

定义 $F=U-TS$，相同的参数项保留为
\[
 \dd F=-S\dd T-p\dd V+Y\dd a.
\]
固定其他变量后，有 $Y=F_a|_T$、$S_a|_T=-Y_T|_a$。进一步改用 $T,a$ 作变量，
\[
 \left(\frac{\partial U}{\partial a}\right)_T
 =Y-T\left(\frac{\partial Y}{\partial T}\right)_a.
\]
它把第二讲的体积公式与表面公式统一起来：取 $a=V,Y=-p$，体积系数为 $-p+Tp_T$；取 $a=\Sigma,Y=\sigma$，面积系数为 $\sigma-T\sigma_T$。两个符号差别源于两种广义功的定义。

\Diagram{fig:generalizedcoordinates}{不同约束方向对应不同的能量系数}{
 \node[draw](u)at(0,0){$U(S,V,a)$};
 \node[draw](s)at(4,1.6){熵方向：$T\dd S$};
 \node[draw](v)at(4,0){体积方向：$-p\dd V$};
 \node[draw](a)at(4,-1.6){参数方向：$Y\dd a$};
 \draw[axis](u)--(s);\draw[axis](u)--(v);\draw[axis](u)--(a);
}
图\ref{fig:generalizedcoordinates}中，每一方向都规定其他独立坐标保持不变。若实际装置同时改变体积和外场参数，两个功项应同时进入第一定律；在实验室能量与旋转参照系的有效能量之间切换时，也要保留变换所引入的角动量项。

\subsection{从微观能量求广义力}
设可分辨的微观状态能量为 $E_i(a)$，没有显式温度依赖。固定 $T,a$ 的配分和为 $Z=\sum_i\ee^{-\beta E_i(a)}$。对参数求导，
\[
 \left(\frac{\partial\ln Z}{\partial a}\right)_T
 =-\beta\left\langle\frac{\partial E_i}{\partial a}\right\rangle,
 \qquad
 \left(\frac{\partial F}{\partial a}\right)_T
 =\left\langle\frac{\partial E_i}{\partial a}\right\rangle.
\]
逐项求导时，$Z_a=\sum_i(-\beta E_{i,a})\ee^{-\beta E_i}$；除以$Z$就得到第一式。再用$F=-\kb T\ln Z$，有$F_a=-\kb T(-\beta\langle E_{i,a}\rangle)=\langle E_{i,a}\rangle$。因此微观参数变化所引起的能量增量，平均后正是宏观广义功。在经典连续相空间中，把求和改成积分也得到同样关系，前提是积分区域和测度的参数依赖已正确处理。

同一配分和还可给出熵。能量$E_i$不显含温度时，固定外参有
\[
 \left(\frac{\partial\ln Z}{\partial T}\right)_a
 =\frac1Z\sum_i\frac{E_i}{\kb T^2}\ee^{-E_i/(\kb T)}
 =\frac{\langle E\rangle}{\kb T^2}.
\]
对$F=-\kb T\ln Z$用乘积法则，得到
\begin{equation}\label{eq:canonicalentropy}
 \begin{aligned}
 S=-\left(\frac{\partial F}{\partial T}\right)_a
 &=\kb\ln Z+\kb T\left(\frac{\partial\ln Z}{\partial T}\right)_a\\
 &=\kb\ln Z+\frac{\langle E\rangle}{T}.
 \end{aligned}
\end{equation}
对$N$个独立经典粒子，$Z_N=Z_1^N/N!$；固定$N$比较状态时，$\ln N!$为常数，在熵差中相消。因而可由单粒子配分积分求出每粒子的熵差，再乘$N$。

若能量零点只是与$a$无关的常数，占有概率、热容与广义力都不变，平均能量只平移同一常数；若减去一个随$a$变化的函数，则$Y$也随之改变。外场参数可变时，不能在不同参数状态各自任选势能零点后再把差值丢掉。例如重力势能应始终相对于同一固定几何高度定义，才能正确计算改变$g$所做的功。

\subsection{广延性、摩尔量与附加界面能}
均匀体相系统在短程相互作用、边界效应可忽略时满足
$U(\lambda S,\lambda V,\lambda N)=\lambda U(S,V,N)$。
对 $\lambda$ 在1处求导，得到 Euler 关系
\[
 U=TS-pV+\mu N.
\]
对右侧三个乘积分别微分，得
\[
 \dd U=T\dd S+S\dd T-p\dd V-V\dd p+\mu\dd N+N\dd\mu.
\]
与$\dd U=T\dd S-p\dd V+\mu\dd N$比较，相同项相消后便有$S\dd T-V\dd p+N\dd\mu=0$。单组分系统的Gibbs自由能为$G=\mu N$；用摩尔量时，化学势按每摩尔定义，便等于摩尔Gibbs自由能。

含显著界面能时，按相同形状把线尺度放大 $\ell$ 倍，体积乘 $\ell^3$，面积却只乘 $\ell^2$。因而把小液滴当作只有体积的均匀广延系统，会漏掉界面项；液滴内外压差与蒸气压的曲率修正都由此产生。远离界面的液相和气相仍按各自的体相摩尔关系计算，界面自由能另行加入，并与面积、体积的几何约束联立。


\begin{example}[总量与摩尔量：实际气体的完整状态函数]\label{ex:vdwmolar}
$\nu$ 摩尔 van der Waals 气体的摩尔体积为 $v=V/\nu$，满足 $(p+a/v^2)(v-b)=RT$。给定常量 $C_{V,m}=c$，推导摩尔内能、摩尔熵和可逆绝热线，再写回总量形式。说明总容积增大与物质的量增大有什么区别。
\Diagram{fig:molarq}{固定温度下分别改变体积和系统尺度}{
 \draw(0,0)rectangle(2,1.4);\node at(1,.7){$\nu,V$};\draw[axis](2.5,.7)--(3.5,.7);\draw(4,0)rectangle(7,1.4);\node at(5.5,.7){$\lambda\nu,\lambda V$};\node at(3.5,-.7){尺度扩大时 $v=V/\nu$ 不变};
}
\end{example}
\begin{solution}
在固定 $\nu$ 下对摩尔体积求偏导，$\partial U_m/\partial v|_T=Tp_T-p=a/v^2$，且 $\partial U_m/\partial T|_v=c$。因此
\[
 U_m-U_{m0}=c(T-T_0)+a(1/v_0-1/v),
\]
\[
 S_m-S_{m0}=c\ln(T/T_0)+R\ln[(v-b)/(v_0-b)].
\]
等熵时 $(v-b)T^{c/R}$ 恒定，等价于 $T(v-b)^{R/c}$ 恒定。图\ref{fig:molarq}保持 $v$ 的尺度变化下，$U$、$S$ 随 $\nu$ 成比例。
从摩尔式写回总量时，先用 $U=\nu U_m$、$S=\nu S_m$，再把所有摩尔体积都换成 $V/\nu$。内能吸引项因而由 $-\nu a/v$ 变成 $-\nu^2a/V$；熵中的体积比则让分子、分母共同含有的 $1/\nu$ 约去。这里的初末态比较始终固定物质的量。随后讨论的 $\lambda$ 倍尺度扩大，是把参考态也复制相同倍数的另一种操作，不能在同一个固定粒子数体积偏导中把 $\nu$ 一起改变。

\Diagram{fig:molara}{摩尔式写回总量后吸引项与排除体积的变化}{
 \node[draw]at(0,0){$a/v^2$};\draw[axis](.8,0)--(2.1,0);\node[draw]at(3.5,0){$\nu^2a/V^2$};
 \node[draw]at(0,-1){$v-b$};\draw[axis](.8,-1)--(2.1,-1);\node[draw]at(3.5,-1){$(V-\nu b)/\nu$};
}
以相同物质的量的参考态比较，
\[
 \Delta U=\nu c\Delta T+\nu^2a(1/V_0-1/V),
\]
\[
 \Delta S=\nu c\ln(T/T_0)+\nu R\ln[(V-\nu b)/(V_0-\nu b)].
\]
图\ref{fig:molara}说明 $a$ 的总容积形式含 $\nu^2$，而排除体积为 $\nu b$。固定 $\nu$ 的等温膨胀会改变摩尔体积和势能；同时放大 $V,\nu$ 则只复制系统，不能把两种变化混用在同一偏导中。
\end{solution}

\subsection{从辐射压强独立推出四次方律}
\label{th:photonlaw}
第二讲由辐射通量的实验温度律写出 $u\propto T^4$。现在反过来，只用各向同性光子的 $p=u/3$、均匀辐射能的广延性 $U=Vu(T)$ 和热力学恒等式。固定温度改变体积时 $U_V=u$，而
\[
 u=-p+Tp_T=-\frac u3+\frac T3\frac{\dd u}{\dd T}.
\]
于是$T\dd u/\dd T=4u$。对$u>0,T>0$分离变量，$\int_{u_0}^{u}\dd u'/u'=4\int_{T_0}^{T}\dd\vartheta/\vartheta$，得到$\ln(u/u_0)=4\ln(T/T_0)$，即$u=aT^4$、$a=u_0/T_0^4$。热力学只能确定幂指数，常数$a$需要辐射谱或实验给出。再由$J_E=cu/4$，得到Stefan--Boltzmann通量律及$\sigma_{\rm SB}=ca/4$。

\begin{example}[辐射腔的可逆绝热关系]\label{ex:radiationadiabat}
腔内只有平衡光子气体，体积由 $V_1$ 缓慢可逆地增至 $V_2$，不与外界换热。求温度、压强、总能量及做功的变化，并与单原子理想气体比较。
\Diagram{fig:radadiabatq}{反射活塞使辐射腔缓慢膨胀}{
 \draw(0,0)rectangle(3.5,1.6);\draw[line width=2pt](3.5,0)--(3.5,1.6);\draw[axis](3.5,.8)--(5,.8);\node at(1.75,.8){$U=aVT^4$};\node at(1.75,-.5){$Q=0$，过程可逆};
}
\end{example}
\begin{solution}
图\ref{fig:radadiabatq}的辐射熵为 $S=4aVT^3/3$。可逆绝热时熵不变，所以
\[
 TV^{1/3}=\text{常量},\quad pV^{4/3}=\text{常量},\quad
 U=aVT^4\propto V^{-1/3}.
\]
因此 $T_2/T_1=(V_1/V_2)^{1/3}$，$p_2/p_1=(V_1/V_2)^{4/3}$，$U_2/U_1=(V_1/V_2)^{1/3}$。
功积分也可逐步与内能差对应。令 $K=p_1V_1^{4/3}$，则
\[
 W=K\int_{V_1}^{V_2}V^{-4/3}\dd V
 =[-3KV^{-1/3}]_{V_1}^{V_2}
 =3K(V_1^{-1/3}-V_2^{-1/3}).
\]
用$U_1=3p_1V_1$改写首项，所得功等于上面的内能减少。能量密度按$T^4$变化，总能量还须乘以当前体积，二者具有不同的体积指数。平衡光子数不受独立的守恒约束；但黑体光子数密度正比于$T^3$，在本题这个保持$VT^3$不变的理想可逆绝热过程中，平均总光子数恰好不变，不能由“光子可产生、湮灭”推断它在每种过程中都会改变。

\Diagram{fig:radadiabata}{等熵膨胀时辐射能降低并转化为功}{
 \draw[axis](0,0)--(5,0)node[right]{$V/V_1$};\draw[axis](0,0)--(0,2.6)node[above]{$U/U_1$};\draw[curve]plot[domain=1:4.6](\x,{2/\x^(1/3)});\node[left]at(0,2){$1$};\node[below]at(1,0){$1$};
}
图\ref{fig:radadiabata}中的能量差为$W=U_1-U_2=U_1[1-(V_1/V_2)^{1/3}]$，与上述压强积分一致。单原子理想气体的绝热指数为$5/3$，平衡辐射则为$4/3$；后者应由辐射状态方程和熵约束推导。
\end{solution}

\section{表面张力系统}
\label{th:surface}
液面面积增大时需要外界做功。以 $\Sigma$ 表示实际界面总面积、$\sigma$ 表示表面张力，在忽略体积功和其他功的模型中
\begin{equation}\label{eq:surfacefund}
 \dd U=T\dd S+\sigma\dd\Sigma.
\end{equation}
这与体积系统的 $-p\dd V$ 符号不同：增大面积需要对系统做正功。表面 Helmholtz 自由能满足 $\dd F=-S\dd T+\sigma\dd\Sigma$，故
\begin{equation}\label{eq:surfacemaxwell}
 \left(\frac{\partial S}{\partial\Sigma}\right)_T
 =-\left(\frac{\partial\sigma}{\partial T}\right)_\Sigma.
\end{equation}
改用 $\Sigma,T$ 作变量，链式法则给
\begin{equation}\label{eq:surfaceenergy}
 \left(\frac{\partial U}{\partial\Sigma}\right)_T
 =\sigma-T\left(\frac{\partial\sigma}{\partial T}\right)_\Sigma,
 \quad
 \dd S=\frac{C_\Sigma}{T}\dd T-\sigma_T\dd\Sigma.
\end{equation}
因此表面内能不一般等于 $\sigma\Sigma$；表面张力是自由能的面积导数。薄膜有两面时 $\Sigma=2A$，单液面才是 $\Sigma=A$。若还计入液体体积热容，必须说明 $C_\Sigma$ 是整个研究系统的定面积热容还是仅表面过剩热容。

\subsection{由几何变化求毛细压强}
\label{th:laplacegeometry}
表面张力给定的是增加实际界面面积所需的等温可逆功。对半径为 $r$ 的球形单界面，
$\Sigma=4\pi r^2$、$V=4\pi r^3/3$，所以
\[
 \frac{\dd\Sigma}{\dd V}
 =\frac{8\pi r}{4\pi r^2}=\frac2r.
\]
虚拟地把界面向外移动 $\dd r$，内外压差提供功
$(p_{\rm in}-p_{\rm out})\dd V$，界面自由能增加
$\sigma\dd\Sigma$。机械平衡要求两者相等，得
\[
 p_{\rm in}-p_{\rm out}=\frac{2\sigma}{r}.
\]
肥皂泡有两层半径近似相同的界面，实际总面积是 $8\pi r^2$，因此压差为 $4\sigma/r$。液滴的液气界面只有一层，不能使用肥皂泡的因子4。

\Diagram{fig:laplacegeometry}{虚位移把球面面积增量与体积增量联系起来}{
 \draw[curve](0,0)circle(1.15);\draw[guide](0,0)circle(1.45);
 \draw[axis](0,0)--(1.15,0)node[midway,below]{$r$};
 \draw[<->](.813,.813)--(1.025,1.025);
 \node[above right]at(1,1){$\dd r$};
 \node[right,align=left]at(2,.4){$\dd\Sigma=8\pi r\,\dd r$\\$\dd V=4\pi r^2\dd r$};
}
图\ref{fig:laplacegeometry}中，虚位移比较相邻几何状态，内外压差与表面张力在这一步作平衡功。圆柱单界面侧面的面积为 $2\pi rL$，体积为 $\pi r^2L$，固定 $L$ 时相同方法给压差 $\sigma/r$。两种几何的差别来自两个主方向中有几个发生弯曲。

\subsection{表面自由能、表面内能与吸热}
对同一面积增量，等温外功为 $\sigma\dd\Sigma$，内能增量为
$(\sigma-T\sigma_T)\dd\Sigma$，可逆吸热为
$-T\sigma_T\dd\Sigma$。三者满足第一定律
$\dd U=\delta Q+\delta W_{\rm on}$。温度升高时表面张力若下降，$\sigma_T<0$，等温拉伸便需要从外界吸热；这些热量进入表面内能，使内能面积系数高于自由能面积系数。

温度改变而面积不变时，热量由定面积热容决定。若把界面和一定量的液体作为整体，$C_\Sigma$ 包含该整体的储热能力；若只讨论表面过剩量，则须把体相的热容另加。选择研究边界后，同一式子中的 $U,S,C_\Sigma$ 均须按这个边界定义。

\begin{example}[液滴合并时的温升与熵变]
\label{ex:dropcoalescence}
$n_d\ge2$ 个相同球形小液滴初始半径为 $r$、温度为 $T_i$，最终合并成一个静止球形液滴。液体不可压缩，忽略蒸发、外界做功和热交换，合并过程产生的宏观动能最终全部耗散。整个液体的定面积热容为常数 $C$，表面张力在所涉温区内为 $\sigma(T)=\sigma_0+\alpha T>0$，其中 $\sigma_0>0$。求终半径、终温和熵变。
\Diagram{fig:dropcoalescenceq}{总体积不变的球形液滴合并}{
 \foreach \x in{0,2.2}\draw[curve](\x,0)circle(.35);
 \node at(1.1,0){$\cdots$};
 \node at(1.1,-.75){$n_d$ 个，半径 $r$};
 \draw[axis](3,0)--(4.1,0);
 \draw[curve](5.2,0)circle(.75);
 \draw[axis](5.2,0)--(5.95,0)node[midway,below]{$R$};
 \node at(5.2,-1.1){$Q=W_{\rm ext}=0$};
}
\end{example}
\begin{solution}
图\ref{fig:dropcoalescenceq}中液体总体积不变，故
$n_dr^3=R^3$，即 $R=n_d^{1/3}r$。初末实际界面面积分别为
\[
 \Sigma_i=4\pi n_dr^2,\qquad
 \Sigma_f=4\pi n_d^{2/3}r^2.
\]
记面积减少量 $D=\Sigma_i-\Sigma_f>0$。由表面内能公式，
$U(T,\Sigma)=CT+\sigma_0\Sigma+\text{常数}$。
整体绝热且无外界功，初末动能均为零，因此
\[
 C(T_f-T_i)=\sigma_0D,\qquad
 T_f=T_i+\frac{\sigma_0D}{C}.
\]
\Diagram{fig:dropcoalescencea}{减少的表面内能转为液体的热运动能}{
 \node[draw](a)at(0,0){$\sigma_0D$};
 \node[draw](b)at(4,0){$C(T_f-T_i)$};
 \draw[axis](a)--(b)node[midway,above]{能量转化};
 \node at(2,-1){外界吸热为零，系统温度仍可升高};
}
图\ref{fig:dropcoalescencea}中的表面能系数是 $\sigma_0=\sigma-T\sigma_T$，而不是初温下的 $\sigma(T_i)$。熵是状态函数，用
$\dd S=C\dd T/T-\alpha\dd\Sigma$ 积分，
\[
 \Delta S=C\ln\frac{T_f}{T_i}+\alpha D.
\]
尽管 $\alpha$ 可为负，总熵在题设有效温区内仍增加。沿相同总能量的连续状态曲线，以已减少面积 $D'$ 为参数，
$T=T_i+\sigma_0D'/C$，于是
\[
 \frac{\dd S}{\dd D'}=\frac{\sigma_0}{T}+\alpha
 =\frac{\sigma(T)}T>0.
\]
积分到 $D$ 即证明 $\Delta S>0$。这个等能状态曲线用于计算和判断端点，真实合并过程中的流动与耗散不必准静态。若算出的终温超出线性表面张力或常热容的适用温区，应重新使用实际函数积分。
\end{solution}


\begin{example}[线性表面张力的内能与熵]\label{ex:linearsurface}
某界面系统在研究温区内满足 $\sigma=\sigma_0+\alpha T$，$\sigma_0,\alpha$ 为常量；定面积热容 $C_\Sigma=C$ 为常量。求相对于 $(\Sigma_0,T_0)$ 的 $U,S$，并检验可积条件。
\Diagram{fig:surfacelinearq}{温度改变时表面张力线性变化}{
 \draw[axis](0,0)--(4.5,0)node[right]{$T$};\draw[axis](0,0)--(0,2.8)node[above]{$\sigma$};\draw[curve](0,2.4)--(4,.8);\node[left]at(0,2.4){$\sigma_0$};\node at(3.3,2.3){示意 $\alpha<0$};
}
\end{example}
\begin{solution}
由图\ref{fig:surfacelinearq}的斜率 $\sigma_T=\alpha$，式\eqref{eq:surfaceenergy}给出 $U_\Sigma|_T=\sigma_0+\alpha T-T\alpha=\sigma_0$。再用 $U_T|_\Sigma=C$，积分得
\[
 U-U_0=\sigma_0(\Sigma-\Sigma_0)+C(T-T_0).
\]
Maxwell 关系给 $S_\Sigma|_T=-\alpha$，定面积方向给 $S_T|_\Sigma=C/T$，所以
\[
 S-S_0=-\alpha(\Sigma-\Sigma_0)+C\ln\frac T{T_0}.
\]
温度项的积分沿定面积路径进行：$\Delta S_T=\int_{T_0}^T C\dd\vartheta/\vartheta$。面积项沿等温路径进行：$\Delta S_\Sigma=\int_{\Sigma_0}^{\Sigma}(-\alpha)\dd\Sigma'$。两者相加时，温度积分产生的对数只含无量纲比值，而面积项的量纲为 $(\mathrm{J\,m^{-2}K^{-1}})\mathrm{m^2}$，也是熵。表面张力虽包含温度项，内能的面积导数中该项被 $T\sigma_T$ 抵消，这正是表面自由能和表面内能需要区分的地方。

\Diagram{fig:surfacelineara}{两条积分方向的熵贡献}{
 \draw[axis](0,0)--(4.6,0)node[right]{$\Sigma$};\draw[axis](0,0)--(0,2.8)node[above]{$T$};\draw[axis](1,.7)--(3.2,.7);\draw[axis](3.2,.7)--(3.2,2.1);\node at(2.1,.3){$-\alpha\Delta\Sigma$};\node[right]at(3.2,1.6){$C\ln(T/T_0)$};
}
图\ref{fig:surfacelineara}的温度项必须为对数，不能写成 $C(T-T_0)$：后者量纲为能量，不是熵。一般可积条件为 $(\partial C_\Sigma/\partial\Sigma)_T=-T\sigma_{TT}$，本题两侧均为零，因此线性表面张力与常热容相容。所有式子只在 $\sigma>0$ 且界面模型成立的温区使用。
\end{solution}

\begin{example}[液膜等温拉伸的热与功]
长 $\ell$ 的活动杆拉伸肥皂膜，膜有两面。温度恒为 $T$，表面张力为 $\sigma_0+\alpha T$。活动杆沿垂直杆方向移动距离 $x$，忽略重力和膜厚变化的体积功。求外力、外界做功、膜吸热和内能变化。
\Diagram{fig:filmq}{两面液膜的活动边界}{
 \draw[curve](0,0)--(0,2)--(4,2);\draw[curve](0,0)--(4,0);\draw[line width=2pt](2.5,0)--(2.5,2);\draw[guide](3.7,0)--(3.7,2);\draw[<->](4.3,0)--(4.3,2)node[midway,right]{$\ell$};\draw[axis](2.5,1)--(3.7,1)node[midway,above]{$x$};
}
\end{example}
\begin{solution}
图\ref{fig:filmq}中单面面积增加 $\ell x$，实际总界面增加 $\Delta\Sigma=2\ell x$。两面各对杆施加 $\sigma\ell$，故缓慢拉伸外力 $F=2\sigma\ell$，对膜做功 $W_{\rm on}=2\sigma\ell x$。
图中杆长 $\ell$ 不变，位移 $\dd x$ 使每一面面积增加 $\ell\dd x$，所以外界微元功为 $F\dd x=\sigma\dd\Sigma$。这同时给出力的二倍系数与功的面积因子，两者应使用同一个总界面面积。等温时膜的熵变化只来自面积项；若液面张力随温度下降，即 $\alpha<0$，新生成界面具有正的熵增，需要从恒温热源吸热。内能式中的 $C(T-T_0)$ 在本题等温条件下为零，但定面积热容本身并不因此为零。

\Diagram{fig:filma}{等温拉伸时的自由能、内能与热量}{
 \node[draw]at(0,0){$\Delta F=\sigma\Delta\Sigma$};\node[draw]at(5,0){$\Delta U=\sigma_0\Delta\Sigma$};\node at(2.5,-.9){$Q=T\Delta S=-\alpha T\Delta\Sigma$};
}
由图\ref{fig:filma}，等温熵增$\Delta S=-\alpha\Delta\Sigma$，吸热$Q=-2\alpha T\ell x$，内能增加$\Delta U=2\sigma_0\ell x$。两种能量输入相加，$Q+W_{\rm on}=2(\sigma_0+\alpha T)\ell x-2\alpha T\ell x=\Delta U$，满足第一定律。
\end{solution}

\begin{example}[界面的绝热温度变化]
线性表面张力、常定面积热容的界面系统，从 $(\Sigma_0,T_0)$ 可逆绝热拉伸到面积 $\Sigma$。求 $T(\Sigma)$、表面张力和外界做功。假设温区内模型始终成立。
\Diagram{fig:surfaceadq}{与外界不换热的可动液膜}{
 \draw(0,0)rectangle(4,1.6);\draw[line width=2pt](2.5,0)--(2.5,1.6);\draw[axis](2.5,.8)--(3.8,.8);\node at(1.2,.8){$Q=0$};\node at(2,-.5){$C_\Sigma=C$};
}
\end{example}
\begin{solution}
图\ref{fig:surfaceadq}的过程可逆，故$\dd S=0$。用已求得的熵微分，$C\dd T/T=\alpha\dd\Sigma$；从$(\Sigma_0,T_0)$积分，有$C[\ln T']_{T_0}^{T}=\alpha[\Sigma']_{\Sigma_0}^{\Sigma}$，即$C\ln(T/T_0)=\alpha(\Sigma-\Sigma_0)$，取指数得
\[
 T=T_0\exp[\alpha(\Sigma-\Sigma_0)/C],\quad
 \sigma=\sigma_0+\alpha T_0\exp[\alpha(\Sigma-\Sigma_0)/C].
\]
外界做功积分中的变温部分以指数函数作原函数，得到
\[
 \int_{\Sigma_0}^{\Sigma}\alpha T_0
 \ee^{\alpha(\Sigma'-\Sigma_0)/C}\dd\Sigma'
 =\left[CT_0\ee^{\alpha(\Sigma'-\Sigma_0)/C}\right]_{\Sigma_0}^{\Sigma}
 =CT_0\bigl[\ee^{\alpha\Delta\Sigma/C}-1\bigr].
\]
它对应内能中的温度变化项，符号在 $\alpha<0$ 的拉伸过程中为负；只要全过程张力为正，总输入功仍为正。膜的降温减小了一部分所需输入，但不取消创建界面的能量代价。

\Diagram{fig:surfaceada}{可逆绝热拉伸时的温度变化}{
 \draw[axis](0,0)--(4.7,0)node[right]{$\Sigma-\Sigma_0$};\draw[axis](0,0)--(0,2.7)node[above]{$T$};\draw[curve]plot[domain=0:4.3](\x,{2.3*exp(-.25*\x)});\node at(2.7,2.4){$\alpha<0$};
}
图\ref{fig:surfaceada}中若 $\alpha<0$，拉伸使温度降低。外界做功等于内能增量，
\[
 W_{\rm on}=\sigma_0\Delta\Sigma+C(T-T_0).
\]
常数张力项给出$\sigma_0\Delta\Sigma$，前面求出的指数项给出$C(T-T_0)$。张力随温度变化，故整个过程的功须沿路径积分。
\end{solution}

\section{相平衡与汽化热}
\label{th:clapeyron}
同一物质的液、气两相共存时，温度和压强相等，两相摩尔 Gibbs 自由能也相等：$g_g=g_l$。沿共存线的相邻状态仍满足此等式。利用 $\dd g=-s\dd T+v\dd p$，
\[
 -(s_g-s_l)\dd T+(v_g-v_l)\dd p=0.
\]
定温定压的可逆汽化每摩尔吸热为 $L_m=T(s_g-s_l)$，从而
\begin{equation}\label{eq:clapeyron}
 \frac{\dd p_s}{\dd T}=\frac{L_m}{T(v_g-v_l)}.
\end{equation}
这就是 Clapeyron 方程。汽化时 $L_m>0$、$v_g>v_l$，共存线斜率为正。熔化时应代入液相减固相的摩尔体积差；冰熔化的体积差为负，熔化潜热为正，故其共存线斜率为负。

还可从跨越两条相邻等温线的微小两相卡诺循环推导。上等温线汽化一摩尔吸热 $L_m$，两条绝热线温差为 $\dd T$；循环面积近似为 $(v_g-v_l)\dd p_s$。可逆效率为 $\dd T/T$，故 $(v_g-v_l)\dd p_s/L_m=\dd T/T$，得到同一方程。用循环法时，面积与吸热必须对应同一物质的量。

\subsection{相平衡条件与稳定相的选择}
\label{th:phasestability}
同温同压的两相处于机械、热学平衡后，仍可能发生物质转移。设从相1向相2可逆转移 $\dd\nu$ 摩尔，总 Gibbs 自由能变化为
\[
 \dd G=(g_2-g_1)\dd\nu.
\]
若 $g_2<g_1$，向相2转移会降低总 $G$，因此相2更稳定；只有 $g_2=g_1$ 时，两相才可在不产生净转移趋势的条件下共存。因而相平衡除了同温、同压，还要求转移物质时两相之间没有自由能降低的方向。

沿共存线改变温度、压强时，两相的 $g$ 继续相等。前文 Clapeyron 方程就是这一相等条件的全微分。斜率的分子是相变熵差，分母是体积差：熵差规定升温倾向哪一相，体积差规定加压倾向哪一相。对通常的液气相变，升温有利于气相，加压有利于液相，两种影响共同形成向右上升的饱和曲线。

\Diagram{fig:phasegibbs}{两相摩尔自由能在共存压强处相等}{
 \draw[axis](0,0)--(5.5,0)node[right]{$p$};
 \draw[axis](0,0)--(0,3.5)node[above]{$g$};
 \draw[curve](.6,1.5)--(4.8,2.2)node[right]{$g_l$};
 \draw[curve](.6,.65)--(4.8,2.95)node[right]{$g_g$};
 \draw[guide](2.831,0)--(2.831,1.872);
 \node[below]at(2.831,0){$p_s(T)$};
 \node at(2.2,3.1){固定温度，斜率为摩尔体积};
}
图\ref{fig:phasegibbs}示意共存点附近的变化。因 $g_p|_T=v$，气相的曲线斜率远大于液相；低于饱和压强时气相自由能较低，高于饱和压强时液相较低。图中的直线只表示局部线性化，跨越较大压强区间时应按各相体积积分。

\subsection{两相平台、杠杆规则与可压缩性}
固定温度下，若每摩尔系统的平均体积 $v$ 位于 $v_l$ 与 $v_g$ 之间，稳定平衡可由液气两相组成。令气相摩尔分数为 $x$，则
\[
 v=(1-x)v_l+xv_g,\qquad
 x=\frac{v-v_l}{v_g-v_l}.
\]
压强保持 $p_s(T)$，体积增大通过增加气相比例实现。此时形式上的等温压缩率很大，反映相比例调整，不能再用单一均匀相的微小密度响应解释。若相变成核受阻，实际系统可能暂时停留在亚稳态，但这需要额外的动力学条件。

\Diagram{fig:leverrule}{共存平台上的总体积由两端相体积加权}{
 \draw[axis](0,0)--(6,0)node[right]{$v$};
 \foreach \x/\t in{1/v_l,3.3/v,5/v_g}{
  \draw(\x,-.13)--(\x,.13);\node[below]at(\x,-.15){$\t$};
 }
 \draw[<->](1,.7)--(3.3,.7)node[midway,above]{$x(v_g-v_l)$};
 \draw[<->](3.3,1.5)--(5,1.5)node[midway,above]{$(1-x)(v_g-v_l)$};
}
图\ref{fig:leverrule}中，靠近液相端点意味着气相比例小。本式使用摩尔体积，权重相应为摩尔分数；采用质量分数或体积分数时，应连同被加权的量一起变换。计算两相热容时，显热项与 $x(T)$ 的变化项都从这一约束出发。

\subsection{等面积条件与自由能相等}
在可解析延拓的状态方程中，沿一条等温线有 $\dd g=v\dd p$。假设两端 $v_l,v_g$ 对应同一压强 $p_s$，积分并分部积分，
\[
 g_g-g_l=\int_{v_l}^{v_g}v\,p'(v)\dd v
 =p_s(v_g-v_l)-\int_{v_l}^{v_g}p(v)\dd v.
\]
要求两端自由能相等，便得到
\[
 \int_{v_l}^{v_g}[p(v)-p_s]\dd v=0.
\]
这正是等面积条件。积分中使用的中间均匀态分支是状态方程的数学延拓，不表示系统在平衡相变时真实经过它们。真实等温压缩沿共存平台改变相比例；曲线在平台上下的有向面积相抵，是两端化学势相等的另一种表达。

\subsection{饱和蒸气压近似的使用层次}
从 Clapeyron 方程到常见的指数蒸气压式，需要依次采用三个近似：液相比气相体积可忽略；气相满足理想气体状态方程；潜热在所讨论温区内近似不变。前两项把斜率化为 $L_m p_s/(RT^2)$，最后一项才允许把 $L_m$ 提到积分号外。

温度接近临界点时，两相体积差和潜热同时趋小，忽略液相体积的近似失效；在高压区，气相也可能明显偏离理想气体。对较宽温区，即使气相足够稀薄，也应保留 $L_m(T)$ 的变化。后面的变潜热模型只改进第三层近似，仍保留前两层，所得幂律乘指数式也应在这一范围内理解。


\begin{example}[饱和蒸气压的指数近似]\label{ex:vaporpressure}
在一段温区内忽略液体摩尔体积，气相视为理想气体，摩尔汽化热 $L_m=L$ 为常量。已知 $p_s(T_0)=p_0$，求 $p_s(T)$ 和压强为 $p$ 时的沸点。
\Diagram{fig:vaporq}{液气共存线上给定一个参考点}{
 \draw[axis](0,0)--(4.8,0)node[right]{$T$};\draw[axis](0,0)--(0,3)node[above]{$p_s$};\draw[curve]plot[domain=.8:4.2](\x,{.12*exp(.65*\x)});\fill(2.3,.535)circle(1.3pt)node[above left]{$T_0,p_0$};
}
\end{example}
\begin{solution}
由图\ref{fig:vaporq}沿共存线的体积差 $v_g-v_l\simeq v_g=RT/p_s$，式\eqref{eq:clapeyron}变为 $\dd p_s/p_s=L\dd T/(RT^2)$。积分时 $\int T^{-2}\dd T=-1/T$，所以
\[
 \ln\frac{p_s(T)}{p_0}=\frac LR\left(\frac1{T_0}-\frac1T\right),\qquad
 p_s(T)=p_0\exp\!\left[\frac LR\left(\frac1{T_0}-\frac1T\right)\right].
\]
求沸点时，先把对数式中的 $1/T_b$ 单独移项，再取倒数，不能直接把 $\ln p$ 看作与温度线性。共存线的局部斜率始终正：增压要求更高的沸点。参考点规定了积分常数，因此指数中的每一项都按相同温标和单位计算；$L$ 是摩尔汽化热，必须与摩尔气体常数 $R$ 配套。若使用单位质量潜热，应同时把 $R$ 换成单位质量气体常数。

\Diagram{fig:vapora}{对数蒸气压与倒温度的线性关系}{
 \draw[axis](0,0)--(4.8,0)node[right]{$1/T$};\draw[axis](0,0)--(0,3)node[above]{$\ln p_s$};\draw[curve](.6,2.5)--(4.2,.5);\node at(3.1,2.5){斜率 $-L/R$};
}
图\ref{fig:vapora}给出实验上常用的线性化。沸点满足 $p_s(T_b)=p$，因此 $T_b=[T_0^{-1}-(R/L)\ln(p/p_0)]^{-1}$。高海拔外压较低，$\ln(p/p_0)<0$，沸点随之降低。靠近临界点时潜热和体积差都变化显著，本近似不能外推。
\end{solution}

\begin{example}[汽化热随温度变化的修正]\label{ex:variablelatent}
仍忽略液体体积，气相理想。气相摩尔定容热容 $c_{v,g}$、液相比热 $c_l$ 均视为常量，已知参考温度 $T_0$ 的汽化热为 $L_0$、蒸气压为 $p_0$。推导 $L_m(T)$ 和修正后的 $p_s(T)$。
\Diagram{fig:latentvariableq}{汽化焓差由两相能量和气相膨胀功组成}{
 \node[draw]at(0,0){液体 $U_{m,l}$};\draw[axis](1.3,0)--(2.8,0)node[midway,above]{$L_m$};\node[draw]at(4.5,0){气体 $U_{m,g}+RT$};
}
\end{example}
\begin{solution}
图\ref{fig:latentvariableq}的定温定压汽化热为焓差 $L_m=U_{m,g}-U_{m,l}+p_s(v_g-v_l)$。在题设近似下 $p_sv_g=RT$，且 $\dd U_{m,g}/\dd T=c_{v,g}$、$\dd U_{m,l}/\dd T\simeq c_l$。定义 $q=c_{v,g}+R-c_l$，则
\[
 L_m(T)=L_0+q(T-T_0)=\Lambda+qT,\quad\Lambda=L_0-qT_0.
\]
代入 $\dd\ln p_s=L_m(T)\dd T/(RT^2)$，分成两项积分，
\[
 \ln\frac{p_s(T)}{p_0}=
 \frac qR\ln\frac T{T_0}+\frac\Lambda R\left(\frac1{T_0}-\frac1T\right).
\]
常数 $\Lambda$ 与参考点汽化热 $L_0$ 不同，它是把线性 $L_m(T)$ 外推到零温时的截距，通常不具备真实零温汽化热的含义。积分时分别有 $\int q\dd T/(RT)=(q/R)\ln T$、$\int\Lambda\dd T/(RT^2)=-\Lambda/(RT)$，所以两个因子中的符号由原函数确定。将最终结果对温度求导，会恢复 $\dd\ln p_s/\dd T=L_m(T)/(RT^2)$，也说明在参考温度附近它与给定的 $L_0$ 连续匹配。

\Diagram{fig:latentvariablea}{修正项分别产生幂律因子和指数因子}{
 \node[draw,align=center]at(0,0){$q/T$\\$(T/T_0)^{q/R}$};\node[draw,align=center]at(5,0){$\Lambda/T^2$\\$\exp[\Lambda(1/T_0-1/T)/R]$};
}
所以图\ref{fig:latentvariablea}的两个因子相乘即 $p_s/p_0$。$q$ 可正可负，取决于两相热容；不能预先把幂指数指定为正。$q=0$ 时回到常潜热公式。这里还忽略了液体焓的压强依赖和气相非理想性，修正温变潜热并不自动消除其他近似误差。
\end{solution}

\begin{example}[熔化曲线的负斜率]
某物质熔化潜热为正，固相摩尔体积 $v_s$ 大于液相 $v_l$。求熔点对压强的局部变化率，并以给定常量近似求压强变化 $\Delta p$ 引起的熔点变化。
\Diagram{fig:meltingq}{固、液体积差决定共存线的斜率}{
 \draw[axis](0,0)--(4.8,0)node[right]{$T$};\draw[axis](0,0)--(0,3)node[above]{$p$};\draw[curve](1,2.6)--(2.8,.4);\node at(1,1.1){固};\node at(3.5,1.8){液};
}
\end{example}
\begin{solution}
图\ref{fig:meltingq}从固到液的 $\Delta v=v_l-v_s<0$，熔化熵增为 $L_m/T>0$。Clapeyron 方程给 $\dd p/\dd T=L_m/[T(v_l-v_s)]<0$。局部取倒数，
\[
 \frac{\dd T_m}{\dd p}=\frac{T_m(v_l-v_s)}{L_m}.
\]
取倒数后，潜热分母为正，斜率的符号由 $v_l-v_s$ 决定。在参考态附近用 $T_{m0}$ 计算斜率，保留的是压强变化的一阶项。压强变化较大时，须沿共存线积分 $\dd T/T=(v_l-v_s)\dd p/L_m$，同时计入各量随状态的变化。图中“加压有利于液相”表示同温下液相的 Gibbs 自由能相对降低，因而共存熔点下降。

\Diagram{fig:meltinga}{加压有利于摩尔体积较小的液相}{
 \node[draw]at(0,0){固体积较大};\draw[axis](1.3,0)--(2.7,0)node[midway,above]{加压};\node[draw]at(4.4,0){液体更有利};\node at(2.2,-.8){共存熔点降低};
}
若变化足够小，可把右侧参数取参考值，$\Delta T_m\simeq T_{m0}(v_l-v_s)\Delta p/L_m$。图\ref{fig:meltinga}从自由能角度说明符号：压强增加时，大体积相的 Gibbs 自由能增加更快。不能仅凭“压强抑制相变”一类措辞判断方向。
\end{solution}

\section{相空间分布与外场中的粒子}
\label{th:phaseprobability}
经典独立粒子的单粒子能量为
\[
 E(\vct r,\vct p)=\frac{p^2}{2m}+E_p(\vct r).
\]
与温度 $T$ 的热源平衡时，相空间密度为 $f(\vct r,\vct p)=Z^{-1}\ee^{-E/(\kb T)}$，概率测度为 $\dd^3r\dd^3p$。若用 $\dd^3r\dd^3p/h^3$ 计数微观状态，则归一化因子也相应改变；固定粒子数的经典位置概率和速度概率不受该常数尺度影响。

由于动能与势能可分离，配分积分也分离：
\[
 Z=\int_{\cal V}\ee^{-E_p/(\kb T)}\dd^3r
 \int_{\mathbb R^3}\ee^{-p^2/(2m\kb T)}\dd^3p.
\]
对位置积分得到 Maxwell 动量分布；对动量积分得到 Boltzmann 位置分布，
\begin{equation}\label{eq:boltzmannposition}
 \rho_P(\vct r)=\frac{\ee^{-E_p(\vct r)/(\kb T)}}
 {\int_{\cal V}\ee^{-E_p(\vct r')/(\kb T)}\dd^3r'}.
\end{equation}
温度相同不意味着粒子密度均匀；势能较低处的粒子更多。反过来，密度随位置改变也不意味着不同高度的平均平动能不同。在等温平衡且动能可分离的条件下，每个位置处的条件速度分布相同。

\subsection{局部密度、位置概率与几何测度}
\label{th:spatialmeasure}
式\eqref{eq:boltzmannposition}给出的是每单位真实空间体积的概率密度。若装置截面随坐标改变，单独观测某一坐标的分布还须计入这一几何因素。例如轴对称圆筒内，体密度为 $n(r)$ 时，半径在 $r$ 与 $r+\dd r$ 间的粒子数为
\[
 \dd N=2\pi rL\,n(r)\dd r.
\]
因而径向概率密度正比于 $r n(r)$，即使 $n(r)$ 为常数，径向样本也更多出现在较大的 $r$。采用 $s=r^2/R_c^2$ 时，$2r\dd r=R_c^2\dd s$，几何因子才被吸收到新变量的微分中。这一几何计数与第一讲速度空间的球壳因子相同：等坐标间隔内所包含的体积会随坐标变化。

\Diagram{fig:spatialannulus}{圆筒横截面的环带测度}{
 \draw[curve](0,0)circle(1.5);
 \fill[gray!20,even odd rule](0,0)circle(1.15)(0,0)circle(.95);
 \draw[curve](0,0)circle(.95);\draw[curve](0,0)circle(1.15);
 \draw[axis](0,0)--(35:.95)node[midway,above]{$r$};
 \draw[<->](315:.95)--(315:1.15);
 \node[right]at(1.05,-.9){$\dd r$};
 \node[right,align=left]at(2.2,.2){体积元 $2\pi rL\dd r$\\径向计数含因子 $r$};
}
图\ref{fig:spatialannulus}中的环带是位置计数区域。重力气柱若截面恒定，其高度概率密度与 $n(h)$ 成正比；变截面容器则应乘截面积 $A(h)$。比较不同高度或半径的“粒子多寡”时，须先说明比较的是等体积、等高度厚度还是等半径厚度。

\subsection{位置涨落怎样决定外场热容}
设外场势能 $\varepsilon(\vct r)$ 不显含温度，容器边界固定。位置配分积分为 $Z_r=\int\ee^{-\beta\varepsilon}\dd^3r$。对 $\beta$ 求导，
\[
 \langle\varepsilon\rangle=-\frac{\partial\ln Z_r}{\partial\beta},
 \qquad
 \frac{\partial\langle\varepsilon\rangle}{\partial\beta}
 =-\bigl(\langle\varepsilon^2\rangle-\langle\varepsilon\rangle^2\bigr).
\]
这里$Z_r'=-\int\varepsilon\ee^{-\beta\varepsilon}\dd^3r$、$Z_r''=\int\varepsilon^2\ee^{-\beta\varepsilon}\dd^3r$，故$\partial_\beta(-Z_r'/Z_r)=-Z_r''/Z_r+(Z_r'/Z_r)^2=-\langle\varepsilon^2\rangle+\langle\varepsilon\rangle^2$。再用$\dd\beta/\dd T=-1/(\kb T^2)$，得单粒子势能对热容的贡献
\[
 \frac{\dd\langle\varepsilon\rangle}{\dd T}
 =\frac{\operatorname{Var}(\varepsilon)}{\kb T^2}\ge0.
\]
因此固定外场中的升温会增加平均势能，增加多少由能量涨落确定。对独立单原子粒子，完整热容为
$C=3N\kb/2+N\operatorname{Var}(\varepsilon)/(\kb T^2)$。
若能量函数本身显含温度，或控制参数在过程中改变，则还需计入相应导数项，不能直接沿用这个固定参数公式。

均匀重力场给 $\varepsilon=mgh$，故气柱的附加热容为
$N(mg)^2\operatorname{Var}(h)/(\kb T^2)$。在有限高度且高温接近均匀分布时，$\operatorname{Var}(h)\to H^2/12$，附加热容趋零；在尺度高度远小于 $H$ 时，高度近似指数分布，其方差为 $(\kb T/(mg))^2$，附加热容趋于 $N\kb$。两个极限也将由后文的有限气柱能量积分得到；方差形式还直接保证附加热容非负。

\subsection{参数变化与可逆绝热降温}
\begin{example}[谐势阱缓慢减弱时的温度与功]
\label{ex:trapadiabat}
$N$ 个经典独立单原子粒子在三维谐势
$\varepsilon=(\kappa_xx^2+\kappa_yy^2+\kappa_zz^2)/2$
中保持热平衡，初温为 $T_1$。在与外界不交换热量的条件下，缓慢改变 $\kappa_x$，使其由 $\kappa_{x1}$ 变为 $\kappa_{x2}>0$，其他两个刚度不变。假定内部碰撞能及时建立平衡，粒子不逃逸，求终温、空间宽度和外界对气体所做的功。
\Diagram{fig:trapadiabatq}{同一谐势方向减弱后的势能曲线}{
 \draw[axis](-2.4,0)--(2.6,0)node[right]{$x$};
 \draw[axis](0,0)--(0,3)node[above]{$\varepsilon_x$};
 \draw[curve]plot[domain=-1.6:1.6](\x,{\x*\x});
 \draw[curve,dashed]plot[domain=-2.2:2.2](\x,{.35*\x*\x});
 \node[right]at(2.2,2.7){初始 $\kappa_{x1}$};
 \node[right]at(2.8,1.2){减弱后 $\kappa_{x2}$};
}
\end{example}
\begin{solution}
图\ref{fig:trapadiabatq}中，势阱减弱改变了可用的位置空间。对动量和位置分别作三个高斯积分，单粒子配分积分的参数依赖为
$Z_1\propto T^3/(\kappa_x\kappa_y\kappa_z)^{1/2}$。其中三个动量高斯积分各给$T^{1/2}$，三个位置高斯积分各给$(T/\kappa_i)^{1/2}$，归一化使用第一讲已证明的高斯积分。按式\eqref{eq:canonicalentropy}，每粒子的熵为$\kb[\ln Z_1+T\partial_T\ln Z_1]=\kb(\ln Z_1+3)$，再加与本题状态无关的固定粒子数常数。比较两态时常数相消，熵差为
\[
 \Delta S=N\kb\left[
 3\ln\frac{T_2}{T_1}
 -\frac12\ln\frac{\kappa_{x2}\kappa_{y2}\kappa_{z2}}
 {\kappa_{x1}\kappa_{y1}\kappa_{z1}}\right].
\]
准静态且内部平衡、无耗散的绝热过程满足 $\Delta S=0$，故记 $a=\kappa_{x2}/\kappa_{x1}$，有
\[
 T_2=T_1a^{1/6}.
\]
每个动能项和势能项各贡献 $\kb T/2$，总能量为 $U=3N\kb T$。因此外界对系统做功
\[
 W_{\rm on}=U_2-U_1=3N\kb T_1(a^{1/6}-1).
\]
\Diagram{fig:trapadiabata}{等熵条件把势阱刚度与温度联系起来}{
 \draw[axis](0,0)--(5,0)node[right]{$\ln(\kappa_x/\kappa_{x1})$};
 \draw[axis](2,-1.1)--(2,1.8)node[above]{$\ln(T/T_1)$};
 \draw[curve](.3,-.283)--(4.6,.433);
 \node[right]at(3.3,1.05){斜率 $1/6$};
 \draw[->](1.4,-.1)--(.7,-.217);
 \node[below]at(1,-.55){减弱势阱};
}
图\ref{fig:trapadiabata}中向左下移动对应 $a<1$：气体降温，并向控制势阱的装置输出功。还可直接计算广义功验证符号，
$\delta W_{\rm on}=\langle\sum_i x_i^2/2\rangle\dd\kappa_x
 =N\kb T\,\dd\kappa_x/(2\kappa_x)$。
令$z=\kappa_x/\kappa_{x1}$，则$\dd\kappa_x/\kappa_x=\dd z/z$，上下限为$1,a$。代入$T=T_1z^{1/6}$，功为$(N\kb T_1/2)\int_1^a z^{-5/6}\dd z=(N\kb T_1/2)[6z^{1/6}]_1^a=3N\kb(T_2-T_1)$，与内能差相同。

位置二阶矩为 $\langle x^2\rangle=\kb T/\kappa_x$，所以
$\sqrt{\langle x^2\rangle_2/\langle x^2\rangle_1}=a^{-5/12}$；
另两方向因刚度固定，宽度之比均为 $a^{1/12}$。
减弱一个方向的束缚会使该方向显著展宽，同时通过碰撞降温使另两方向收窄。若三个刚度一同乘 $a$，等熵式改给 $T_2/T_1=a^{1/2}$，可见不同参数操作不能共用同一指数。
\end{solution}


\subsection{等温大气的两种推导}
\label{th:barometric}
设重力加速度$g$恒定，气体摩尔质量为$M$、温度为常量$T$。高度$h$到$h+\dd h$的薄层受力平衡，$p(h)S-p(h+\dd h)S-\rho Sg\dd h=0$，所以$\dd p/\dd h=-\rho g$。理想气体$\rho=Mp/(RT)$，分离变量并从底面积分至高度$h$，有
\[
 \int_{p(0)}^{p(h)}\frac{\dd p'}{p'}
 =-\frac{Mg}{RT}\int_0^h\dd h',\qquad
 \ln\frac{p(h)}{p(0)}=-\frac{Mgh}{RT}.
\]
取指数，再用$n=p/(\kb T)$，得到
\begin{equation}\label{eq:barometric}
 p(h)=p(0)\ee^{-Mgh/(RT)},\qquad
 n(h)=n(0)\ee^{-mgh/(\kb T)}.
\end{equation}
用式\eqref{eq:boltzmannposition}，单粒子势能为 $mgh$，立即得到同一密度指数。两条路线分别从宏观受力和微观概率进入，指数尺度一致，因为 $M/R=m/\kb$。尺度高度为 $H_T=RT/(Mg)$。地球实际大气温度随高度变化，不能在任意高度直接套用等温指数律。

\begin{example}[气柱压强与总重量]\label{ex:atmosphere}
截面积为 $S$、高 $H$ 的竖直刚性柱内，封闭 $N$ 个质量为 $m$ 的理想气体粒子，处于温度 $T$ 的平衡。求底面与顶面压强，并证明其差产生的合力等于气体重量。
\Diagram{fig:columnpressureq}{上下两端的压强不同}{
 \draw(0,0)rectangle(2,3);\draw[axis](2.6,0)--(2.6,3.5)node[above]{$h$};\node[left]at(0,0){$p_b$};\node[left]at(0,3){$p_t$};\node at(1,1.5){$N,T$};\draw[<->](3.3,0)--(3.3,3)node[midway,right]{$H$};
}
\end{example}
\begin{solution}
取 $a_g=mg/(\kb T)$、$x=a_gH$，密度为 $n(h)=n_b\ee^{-a_gh}$。图\ref{fig:columnpressureq}的固定粒子数要求
\[
 N=S\int_0^Hn(h)\dd h=Sn_b\frac{1-\ee^{-x}}{a_g},
 \quad n_b=\frac{Na_g}{S(1-\ee^{-x})}.
\]
因此
\[
 p_b=\frac{Nmg}{S(1-\ee^{-x})},\qquad
 p_t=p_b\ee^{-x}.
\]
这里高度积分为$\int_0^H\ee^{-a_gh}\dd h=[-\ee^{-a_gh}/a_g]_0^H=(1-\ee^{-x})/a_g$。底部密度由总粒子数决定，并非外界气源预先固定的密度。乘上$\kb T$才得到压强，$a_g\kb T=mg$使显式分子中的温度因子约去，但温度仍保留在指数里。

\Diagram{fig:columnpressurea}{气体整体的竖直受力平衡}{
 \draw(0,0)rectangle(2,2.5);\draw[axis](1,0)--(1,1)node[right]{$Sp_b$};\draw[axis](1,2.5)--(1,1.8)node[right]{$Sp_t$};\draw[axis](2.6,1.6)--(2.6,.5)node[right]{$Nmg$};
}
图\ref{fig:columnpressurea}中的气体所受底面向上力减去顶面向下力为 $S(p_b-p_t)=Nmg$，恰好平衡总重力。升温时底面密度归一化常数也会改变，不能把 $n_b$ 当与温度无关的固定量。
\end{solution}

\begin{example}[有限高度气柱的总热容]\label{ex:columnheat}
上述刚性封闭气柱中装有 $\nu$ 摩尔单原子理想气体，摩尔质量 $M$。缓慢供热使每一状态都处于同一均匀温度的热平衡。求重力势能 $E_p(T)$ 和容器的热容 $C=\delta Q/\dd T$，忽略容器自身热容。
\Diagram{fig:columnheatq}{升温除了增大热运动，还改变高度分布}{
 \draw(0,0)rectangle(2,3);\foreach \y in{.25,.55,.85,1.25,1.8,2.5}{\fill(.8,\y)circle(1.2pt);}\draw[axis](-1,1.3)--(0,1.3)node[midway,above]{$Q$};\node[right]at(2.2,1.5){$E_p=\int_0^H Mgh\,\dd\nu$};
}
\end{example}
\begin{solution}
图\ref{fig:columnheatq}的气体密度从底部向上衰减，令 $x=MgH/(RT)$、$y=h/H$。按上一例归一化，单粒子高度的无量纲密度为 $w(y)=x\ee^{-xy}/(1-\ee^{-x})$。于是
\[
 E_p=\nu MgH\frac{x}{1-\ee^{-x}}\int_0^1y\ee^{-xy}\dd y.
\]
分部积分取 $u=y$、$\dd v=\ee^{-xy}\dd y$，
\[
 \int_0^1y\ee^{-xy}\dd y
 =\left[-\frac{y\ee^{-xy}}x\right]_0^1
   +\frac1x\left[-\frac{\ee^{-xy}}x\right]_0^1
 =-\frac{\ee^{-x}}x+\frac{1-\ee^{-x}}{x^2}.
\]
所以
\begin{equation}\label{eq:columnpotential}
 E_p=\nu MgH\left[\frac1x-\frac1{\ee^x-1}\right]
 =\nu RT-\frac{\nu MgH}{\ee^x-1}.
\end{equation}
对势能式求温度导数时，第二项的全部温度依赖都藏在 $x=MgH/(RT)$ 中。链式法则给
\[
 \frac{\dd}{\dd T}\frac1{\ee^x-1}
 =-\frac{\ee^x}{(\ee^x-1)^2}\left(-\frac xT\right)
 =\frac{x\ee^x}{T(\ee^x-1)^2}.
\]
再用 $MgH=xRT$，把这一项的系数写为无量纲 $x^2$，便得到热容中的修正分式。系统能量包含外场势能，所以固定几何体积下的供热量不只是平动内能变化；若把势能留在系统之外，则还应等价地计入重力对粒子所做的功。

\Diagram{fig:columnheata}{刚性容器的供热分给平动能与重力势能}{
 \node[draw]at(0,0){$\delta Q$};\draw[axis](.8,.2)--(2.3,1);\draw[axis](.8,-.2)--(2.3,-1);\node[right]at(2.3,1){$\dd U_{\rm tr}=3\nu R\dd T/2$};\node[right]at(2.3,-1){$\dd E_p$};
}
外壁刚性，无边界体积功。按图\ref{fig:columnheata}将重力势能计入系统，有$\delta Q=\dd(3\nu RT/2+E_p)$。用前面求出的势能导数，再加上平动热容，得到
\begin{equation}\label{eq:columncapacity}
 C=\nu R\left[\frac52-\frac{x^2\ee^x}{(\ee^x-1)^2}\right].
\end{equation}
式中第一项 $3\nu R/2$ 是局部平动热容，额外部分来自高度分布重排。气体内部没有活塞运动，仍可有超过定容平动热容的供热需求，因为外场势能随温度改变。
\end{solution}

\begin{example}[气柱热容的两端极限]\label{ex:columnlimits}
分析式\eqref{eq:columncapacity}在 $x\ll1$、$x\gg1$ 时的行为；求平均高度的相应极限，并说明先令 $H\to\infty$ 与先令 $T\to\infty$ 为什么不同。
\Diagram{fig:columnlimitq}{尺度高度相对于容器高度的两种情形}{
 \draw(0,0)rectangle(1.8,3);\foreach \y in{.3,.8,1.3,1.8,2.3,2.8}\fill(.9,\y)circle(1.3pt);\node at(.9,-.5){$x\ll1$};
 \draw(3.4,0)rectangle(5.2,3);\foreach \y in{.15,.3,.5,.75,1.1,1.7}\fill(4.3,\y)circle(1.3pt);\node at(4.3,-.5){$x\gg1$};
}
\end{example}
\begin{solution}
图\ref{fig:columnlimitq}左图$H\ll H_T$，密度近乎均匀。先从指数的Taylor展开求倒数。记$F(x)=1/(\ee^x-1)$，写成
\[
 \begin{aligned}
 \ee^x-1&=x\left(1+\frac x2+\frac{x^2}6+\frac{x^3}{24}
 +\frac{x^4}{120}+\cdots\right),\\
 F(x)&=\frac1x(1+a_1x+a_2x^2+a_3x^3+a_4x^4+\cdots).
 \end{aligned}
\]
两式相乘须为1，各次幂系数依次满足
\[
 \begin{aligned}
 a_1+\tfrac12&=0,\\
 a_2+\tfrac12a_1+\tfrac16&=0,\\
 a_3+\tfrac12a_2+\tfrac16a_1+\tfrac1{24}&=0,\\
 a_4+\tfrac12a_3+\tfrac16a_2+\tfrac1{24}a_1+\tfrac1{120}&=0.
 \end{aligned}
\]
逐式回代得$a_1=-1/2$、$a_2=1/12$、$a_3=0$、$a_4=-1/720$。又有$F(-x)=-1-F(x)$，故$F(x)+1/2$的Laurent展开只有奇次幂，得到
\[
 F(x)=\frac1x-\frac12+\frac x{12}-\frac{x^3}{720}+O(x^5).
\]
式\eqref{eq:columnpotential}中的$1/x$先相消，平均高度为
\[
 \langle h\rangle=H\left(\frac1x-F(x)\right)
 =H\left(\frac12-\frac x{12}+\frac{x^3}{720}+O(x^5)\right).
\]
因此首两项为$H/2-MgH^2/(12RT)$。热容须对能量求导，不能把高度展开直接当成热容展开。对$F$逐项求导，
\[
 F'(x)=-\frac1{x^2}+\frac1{12}-\frac{x^2}{240}+O(x^4).
\]
由$E_p=\nu RT-\nu MgHF(x)$及$\dd x/\dd T=-x/T$，有
\[
 \frac{\dd E_p}{\dd T}=\nu R+\frac{\nu MgHx}{T}F'(x)
 =\nu R[1+x^2F'(x)].
\]
加上平动热容，得到
\[
 C=\nu R\left[\frac52+x^2F'(x)\right]
 =\nu R\left[\frac32+\frac{x^2}{12}-\frac{x^4}{240}+O(x^6)\right].
\]
有限顶壁使平均势能在高温下趋于$\nu MgH/2$，附加热容从正值趋零。右图$H\gg H_T$时，$F(x)=\ee^{-x}[1+O(\ee^{-x})]$，顶壁影响指数小，$E_p\simeq\nu RT$、$\langle h\rangle\simeq RT/(Mg)$，故$C\simeq5\nu R/2$。

\Diagram{fig:columnlimita}{气柱热容随无量纲高度单调升高}{
 \draw[axis](0,0)--(5.5,0)node[right]{$x$};\draw[axis](0,0)--(0,3.1)node[above]{$C/(\nu R)$};
 \draw[curve]plot[domain=.05:6,samples=80]({.8*\x},{2.5-\x*\x*exp(-\x)/(1-exp(-\x))^2});\draw[guide](0,2.5)--(5,2.5);\node[left]at(0,1.5){$3/2$};\node[left]at(0,2.5){$5/2$};
}
图\ref{fig:columnlimita}的单调性也可由公式确定。将修正分式写成$h(x)^2$，其中$h(x)=x/[2\sinh(x/2)]$。令$y=x/2>0$，商法则给
\[
 h'(x)=\frac{\sinh y-y\cosh y}{2\sinh^2y}.
\]
分子反号后的函数$f(y)=y\cosh y-\sinh y$满足$f(0)=0$、$f'(y)=y\sinh y>0$，故$h'<0$，于是$\dd[C/(\nu R)]/\dd x=-2hh'>0$，曲线从$3/2$单调升至$5/2$。

固定有限$H$先升到高温，粒子高度受顶壁限制，额外热容趋零；先把顶壁移到无穷远，平均高度可以随温度不断增长，始终保留$\nu R$的重力热容。两种操作改变了不同的几何约束，极限不交换。
\end{solution}

\section{界面、相变与外场的综合模型}
液膜拉伸时，外界做功先与表面自由能相联系；计算供热还需表面内能。外场中的气体则同时涉及热运动、势能与外界机械功。以下各题按图中的系统边界分别列出这些能量项。

\begin{example}[非线性表面张力与相容热容]
界面满足 $\sigma=\sigma_0-a_sT^2$，$a_s>0$，热容在零面积参考处为常量 $C_0$。忽略其他广义功。求相容的 $C_\Sigma(\Sigma,T)$、$U$ 与 $S$；说明能否同时规定定面积热容处处等于 $C_0$。
\Diagram{fig:nonlinearsurfaceq}{表面张力曲率影响面积方向的热容变化}{
 \draw[axis](0,0)--(4.6,0)node[right]{$T$};\draw[axis](0,0)--(0,2.8)node[above]{$\sigma$};\draw[curve]plot[domain=0:4](\x,{2.5-.1*\x*\x});\node at(3,2.6){$\sigma_{TT}=-2a_s$};
}
\end{example}
\begin{solution}
图\ref{fig:nonlinearsurfaceq}的 $U_\Sigma=\sigma-T\sigma_T=\sigma_0+a_sT^2$，先对面积积分，$U=(\sigma_0+a_sT^2)\Sigma+\Phi(T)$。在零面积参考处 $U_T=C_0$，故 $\Phi=C_0T+\text{常量}$。因此
\[
 C_\Sigma=C_0+2a_sT\Sigma,
 \qquad U=\sigma_0\Sigma+a_s\Sigma T^2+C_0T+\text{常量}.
\]
Maxwell 关系给 $S_\Sigma=-\sigma_T=2a_sT$，积分并匹配 $S_T=C_\Sigma/T$，得到
\[
 S=2a_s\Sigma T+C_0\ln T+\text{常量}.
\]
给定的 $C_0$ 用来确定温度积分函数 $\Phi(T)$，并非规定所有面积上的热容都相同。对 $U$ 中的 $a_s\Sigma T^2$ 求温度导数，产生不可略去的 $2a_s\Sigma T$。再用所得熵求回基本关系，温度方向有 $T S_T=C_0+2a_s\Sigma T$，面积方向有 $T S_\Sigma+\sigma=\sigma_0+a_sT^2$，分别与内能的两个偏导匹配。两个方向同时匹配，才说明 $U,S$ 属于同一个相容模型。

\Diagram{fig:nonlinearsurfacea}{两种顺序求混合偏导必须一致}{
 \node[draw]at(0,0){$\partial_\Sigma C_\Sigma$};\node at(2.3,0){$=$};\node[draw]at(4.6,0){$-T\sigma_{TT}=2a_sT$};
}
图\ref{fig:nonlinearsurfacea}的等式要求热容随面积变化。若处处规定 $C_\Sigma=C_0$，左侧为零而右侧非零，模型便不相容。线性 $\sigma(T)$ 的二阶温度导数为零，才允许采用面积无关的热容。零面积参考值用于指定积分函数，物理应用范围仍须满足界面可作连续描述且 $\sigma>0$。
\end{solution}

\begin{example}[封闭肥皂泡的等温膨胀]
肥皂泡中有 $\nu$ 摩尔理想气体，温度保持 $T$。薄膜张力 $\sigma=\sigma_0+\alpha T$，两面面积近似相同，泡半径由 $r_1$ 缓慢变为 $r_2$。外压受控制，忽略膜厚、重力和膜体积功。求平衡外压、气体做功、对外负载做功及气体和膜整体吸热。
\Diagram{fig:bubbleq}{泡内外压强差由两层界面提供}{
 \draw[curve](0,0)circle(1.5);\draw[curve](0,0)circle(1.57);\draw[axis](0,0)--(1.5,0)node[midway,above]{$r$};\node at(0,-.65){$p_{\rm in}$};\node at(2.8,.7){$p_{\rm ext}$};\node at(3,-.2){$p_{\rm in}-p_{\rm ext}=4\sigma/r$};
}
\end{example}
\begin{solution}
图\ref{fig:bubbleq}中 $V=4\pi r^3/3$、实际界面面积 $\Sigma=8\pi r^2$，理想气体压强 $p_{\rm in}=\nu RT/V$。机械平衡要求
\[
 p_{\rm ext}(r)=\frac{3\nu RT}{4\pi r^3}-\frac{4\sigma}{r}.
\]
路径只能在外压条件可实现的半径区间运行。气体等温做功为
\[
 W_g=\nu RT\int_{V_1}^{V_2}\frac{\dd V}{V}
 =\nu RT[\ln V]_{V_1}^{V_2}=3\nu RT\ln\frac{r_2}{r_1}.
\]
其中一部分增加膜的自由能，另一部分付给外压负载：
\[
 W_{\rm ext}=\int p_{\rm ext}\dd V
 =W_g-\sigma\Delta\Sigma,
 \quad \Delta\Sigma=8\pi(r_2^2-r_1^2).
\]
膜所需的功也可从气体内外压差直接积分。球体体积微元为 $\dd V=4\pi r^2\dd r$，因此
\[
 \int_{r_1}^{r_2}(p_{\rm in}-p_{\rm ext})\dd V
 =\int_{r_1}^{r_2}16\pi\sigma r\dd r
 =[8\pi\sigma r^2]_{r_1}^{r_2}=8\pi\sigma(r_2^2-r_1^2).
\]
这与两面总面积的自由能增量相同。气体把这部分能量交给膜，膜又与恒温热源交换热；对气体和膜的整体列式时，这笔内部功抵消，外界负载只收到由外压积分计算的那一部分。

\Diagram{fig:bubblea}{以气体和膜为整体时，膜功属于内部能量交换}{
 \node[draw]at(0,0){气体 $W_g$};\draw[axis](1.2,.2)--(2.8,1);\draw[axis](1.2,-.2)--(2.8,-1);\node[right]at(2.8,1){膜自由能增加 $\sigma\Delta\Sigma$};\node[right]at(2.8,-1){外界负载 $W_{\rm ext}$};
}
气体内能不变，膜内能增量为 $\sigma_0\Delta\Sigma$。按图\ref{fig:bubblea}选择整体边界，只有 $W_{\rm ext}$ 是对外功，故
\[
 Q=\Delta U_{\rm total}+W_{\rm ext}
 =W_g-\alpha T\Delta\Sigma.
\]
也可由整体熵变 $\Delta S=\nu R\ln(V_2/V_1)-\alpha\Delta\Sigma$ 得 $Q=T\Delta S$。若把气体做功和膜功都再次计为整体对外功，会重复计算内部交换。
\end{solution}

\begin{example}[定容两相系统的表观热容]
$\nu$ 摩尔物质在刚性容器中处于液气共存，摩尔平均体积为 $\bar v=V/\nu$。气相理想，液体体积可忽略，气相摩尔定容热容为 $c_{v,g}$，液相比热为 $c_l$。沿共存线缓慢加热，摩尔汽化热为 $L_m(T)$。求气相摩尔分数和系统定容供热率 $C=\delta Q/\dd T$。
\Diagram{fig:twophasecapacityq}{容器定容，液气比例随温度改变}{
 \draw(0,0)rectangle(3,2.6);\fill[gray!20](0,0)rectangle(3,.8);\draw[curve](0,.8)--(3,.8);\node at(1.5,.4){液相};\node at(1.5,1.7){气相 $x\nu$};\draw[axis](-1,1)--(0,1)node[midway,above]{$Q$};
}
\end{example}
\begin{solution}
图\ref{fig:twophasecapacityq}的气体占据几乎全部体积，故 $V\simeq x\nu RT/p_s(T)$，于是
\[
 x(T)=\frac{\bar v p_s(T)}{RT},\qquad
 \frac{x'}x=\frac{L_m(T)}{RT^2}-\frac1T.
\]
只在 $0<x<1$ 且两相近似成立的区间使用。总内能为 $U=\nu[(1-x)u_l+xu_g]$，沿共存线求导要保留组分比例的变化：
\[
 \frac{\dd U}{\dd T}=\nu[(1-x)c_l+xc_{v,g}+(u_g-u_l)x'].
\]
由汽化焓差 $L_m=u_g-u_l+RT$，得到
\[
 C=\nu\left[(1-x)c_l+xc_{v,g}
 +(L_m-RT)x\left(\frac{L_m}{RT^2}-\frac1T\right)\right].
\]
由 $x'/x=(L_m-RT)/(RT^2)$，相比例项还可化成
\[
 (L_m-RT)x'=\frac{x(L_m-RT)^2}{RT^2}\ge0.
\]
因此在本题近似下，这一额外储热能力非负。它与气相比例改变的方向无关：升温可使 $x$ 增大，也可在某些参数范围使其减小，相变对应的内能变化会同时换号。容器一旦完全气化或液化，$x$ 不能再按共存线公式继续变化，需改用对应单相的状态方程和热容。

\Diagram{fig:twophasecapacitya}{供热同时改变各相温度和相的数量}{
 \node[draw,align=center]at(0,0){显热项\\$(1-x)c_l+xc_{v,g}$};\node at(2.8,0){$+$};\node[draw,align=center]at(5.1,0){相比例项\\$(L_m-RT)x'$};
}
图\ref{fig:twophasecapacitya}中的相比例项可能很大。容器定容只使外壁体积功为零，内部相变仍会贡献热容。潜热 $L_m$ 是两相焓差，扣除气相形成所需的 $RT$ 后，才得到内能差 $L_m-RT$，故应由它乘以 $x'$。更精确的模型还须保留液体体积及各相沿共存线的压强效应。
\end{solution}

\begin{example}[临界点与等面积条件]
对一摩尔 van der Waals 气体求临界参数。低于临界温度时，说明为什么不能把等温线中 $p_v|_T>0$ 的部分作为稳定平衡路径，并从两相 Gibbs 自由能相等导出等面积条件。
\Diagram{fig:criticalvdwq}{低于临界温度的等温线与共存平台}{
 \draw[axis](0,0)--(5.5,0)node[right]{$v/v_c$};\draw[axis](0,0)--(0,3)node[above]{$p/p_c$};
 \draw[curve]plot[domain=.55:4,samples=120]({1.2*\x},{2*(7.2/(3*\x-1)-3/(\x*\x))});
 \draw[guide](.7241,1.294)--(2.8186,1.294);\node at(3.7,2.5){$T=0.9T_c$};\node[below]at(.7241,0){$v_l/v_c$};\node[below]at(2.8186,0){$v_g/v_c$};
}
\end{example}
\begin{solution}
临界等温线在拐点有 $p_v=0$、$p_{vv}=0$：
\[
 -\frac{RT_c}{(v_c-b)^2}+\frac{2a}{v_c^3}=0,\qquad
 \frac{2RT_c}{(v_c-b)^3}-\frac{6a}{v_c^4}=0.
\]
两式相除消去 $T_c$，得 $2/(v_c-b)=3/v_c$，所以
\[
 v_c=3b,\quad T_c=\frac{8a}{27Rb},\quad p_c=\frac a{27b^2}.
\]
图\ref{fig:criticalvdwq}中向右上升的曲线段满足等温压缩率为负，不能构成稳定均匀相。平衡时由液气两相共存平台替代。
等面积条件不是对实际两相过程使用不稳定均匀状态，而是借数学状态方程比较两端纯相的自由能。实际共存状态由相比例改变连接两端，压强保持 $p_s$。对于平均摩尔体积 $\bar v$ 介于两端之间，气相摩尔分数由 $\bar v=(1-x)v_l+xv_g$ 给出，故 $x=(\bar v-v_l)/(v_g-v_l)$。平台内部的不同横坐标代表不同相比例，而不是均匀气体沿回折线逐点移动。

\Diagram{fig:criticalvdwa}{平台上下的有向面积必须相抵}{
 \draw[axis](0,0)--(5.4,0)node[right]{$v/v_c$};\draw[axis](0,0)--(0,2.6)node[above]{$p/p_c$};
 \draw[curve]plot[domain=.6034:2.34884,samples=100]({1.8*\x},{2.5*(7.2/(3*\x-1)-3/(\x*\x))});
 \draw[curve](1.0861,1.6175)--(4.2279,1.6175);\node at(1.7,.8){负面积};\node at(3.2,2){正面积};
}
在定温下 $\dd g=v\dd p$，沿数学等温线从 $v_l$ 积分到 $v_g$，分部积分得到
\[
 g_g-g_l=[vp]_{v_l}^{v_g}-\int_{v_l}^{v_g}p(v,T)\dd v.
\]
两端共存压强均为 $p_s$，且 $g_g=g_l$，所以
\[
 \int_{v_l}^{v_g}[p(v,T)-p_s]\dd v=0.
\]
这就是图\ref{fig:criticalvdwa}的等面积条件。与两端方程 $p(v_l,T)=p(v_g,T)=p_s$ 联立，可确定液、气摩尔体积与饱和压强。图中取 $T/T_c=0.9$，对应 $v_l/v_c\approx0.60340$、$v_g/v_c\approx2.34884$、$p_s/p_c\approx0.64700$，平台按这些数值绘制。
\end{solution}

\begin{example}[有限气柱的等温升顶做功]
气柱截面积 $S$ 固定，内有 $N$ 个理想气体粒子，处于均匀温度 $T$ 的重力平衡。顶壁可缓慢从高度 $H_1$ 升至 $H_2$，粒子数固定，过程可逆等温。求顶壁做功和吸热；讨论 $H_2\to\infty$ 时的功是否发散。
\Diagram{fig:raisetopq}{升高顶壁时只受顶端局部压强做功}{
 \draw(0,0)--(0,3.4);\draw(2.2,0)--(2.2,3.4);\draw(0,0)--(2.2,0);\draw[line width=1.4pt](0,2)--(2.2,2);\draw[guide](0,3.1)--(2.2,3.1);\draw[axis](1.1,2.1)--(1.1,3);\node[right]at(2.4,2){$H_1$};\node[right]at(2.4,3.1){$H_2$};
}
\end{example}
\begin{solution}
由有限气柱归一化，图\ref{fig:raisetopq}的顶端压强为$p_t=Nmg/[S(\ee^{a_gH}-1)]$，$a_g=mg/(\kb T)$。气体向外做功只在顶壁。令$x=a_gH$，则$\dd H=\dd x/a_g$，上下限由$H_1,H_2$变为$a_gH_1,a_gH_2$，所以
\begin{align*}
 W&=\int_{H_1}^{H_2}Sp_t\dd H
 =N\kb T\int_{a_gH_1}^{a_gH_2}\frac{\dd x}{\ee^x-1}\\
 &=N\kb T\ln\frac{1-\ee^{-a_gH_2}}{1-\ee^{-a_gH_1}}.
\end{align*}
最后一步利用 $\dd\ln(1-\ee^{-x})/\dd x=1/(\ee^x-1)$。温度不变，平动内能不变，重力势能却改变，
\[
 \Delta E_p=Nmg\left[\frac{H_1}{\ee^{a_gH_1}-1}
 -\frac{H_2}{\ee^{a_gH_2}-1}\right].
\]
吸热中的重力势能项取决于平均高度，顶壁功则由顶端的局部压强决定。也可用配分积分求熵：固定温度下，位置积分 $Z_h=(1-\ee^{-a_gH})/a_g$ 随顶壁高度改变，熵差为 $N\kb\ln(Z_{h,2}/Z_{h,1})+\Delta E_p/T$。乘以温度，第一项给出顶壁功，第二项给出势能增量，与第一定律一致。位置积分在无限高时收敛，功也趋于有限值。

\Diagram{fig:raisetopa}{等温供热同时支付重力势能和顶壁功}{
 \node[draw]at(0,0){$Q$};\draw[axis](.7,.2)--(2.3,1);\draw[axis](.7,-.2)--(2.3,-1);\node[right]at(2.3,1){$\Delta E_p$};\node[right]at(2.3,-1){$W$};
}
图\ref{fig:raisetopa}给 $Q=\Delta E_p+W$。$H_2\to\infty$ 时功趋于有限值 $-N\kb T\ln(1-\ee^{-a_gH_1})$，因为高处密度指数小，后续顶壁压强很快趋零。忽略重力时的 $N\kb T\ln(H_2/H_1)$ 发散结论不能用于这个极限；应先比较尺度高度与容器高度。
\end{solution}

\begin{example}[恒定角速度下的旋转气体热容]
圆柱容器在电机维持的恒定角速度 $\Omega$ 下转动，半径 $R_c$，内有 $N$ 个质量为 $m$ 的单原子理想气体粒子。忽略重力，气体相对容器始终处于温度 $T$ 的平衡。求平均转动惯量，正确计入电机功后求供热热容。
\Diagram{fig:rotcapacityq}{固定角速度下的供热与有符号机械功}{
 \draw[curve](0,0)circle(1.6);\draw[->](77:1.9)arc(77:170:1.9);\node at(-1.5,2.2){$\Omega$ 固定};
 \draw[<->](2.2,0)--(3.6,0)node[right]{电机 $\Omega\dd L_z$};\draw[axis](-3,0)--(-1.7,0)node[midway,above]{$Q$};
}
\end{example}
\begin{solution}
令 $a=m\Omega^2/(2\kb T)$、$b_\Omega=aR_c^2$。径向密度 $\propto\ee^{ar^2}$，体积元含 $r\dd r$。对归一化积分求参数导数，
\[
 \langle r^2\rangle=\frac{\dd}{\dd a}\ln\frac{\ee^{aR_c^2}-1}{a}
 =\frac{R_c^2}{1-\ee^{-b_\Omega}}-\frac1a.
\]
转动惯量 $I=Nm\langle r^2\rangle$，角动量 $L_z=I\Omega$。实验室系总能量为平动热能加整体旋转能，$E=3N\kb T/2+I\Omega^2/2$。图\ref{fig:rotcapacityq}中电机为保持 $\Omega$ 还要做功 $\delta W_{\rm on}=\Omega\dd L_z=\Omega^2\dd I$，故
\[
 \delta Q=\dd E-\Omega\dd L_z
 =\frac32N\kb\dd T-\frac12\Omega^2\dd I.
\]
将惯量写为$I=NmR_c^2/(1-\ee^{-b_\Omega})-2N\kb T/\Omega^2$。因$\dd b_\Omega/\dd T=-b_\Omega/T$，第一项求导为$NmR_c^2b_\Omega\ee^{-b_\Omega}/[T(1-\ee^{-b_\Omega})^2]$，第二项为$-2N\kb/\Omega^2$。再用$mR_c^2=2\kb Tb_\Omega/\Omega^2$，合并得
\[
 \frac{\dd I}{\dd T}=-\frac{2N\kb}{\Omega^2}
 \left[1-\frac{b_\Omega^2\ee^{-b_\Omega}}
 {(1-\ee^{-b_\Omega})^2}\right].
\]
括号非负，故升温时径向密度趋于均匀，平均惯量减小。角速度固定，电机须取走一部分角动量，电机对气体做的功因而为负。把这一有符号的机械功从实验室系能量增量中减去，才得到供热量。

\Diagram{fig:rotcapacitya}{升温使径向位置分布变均匀，同时与电机交换功}{
 \draw[axis](0,0)--(4.8,0)node[right]{$s=r^2/R_c^2$};\draw[axis](0,0)--(0,4.7)node[above]{$q(s)$};
 \draw[curve]plot[domain=0:1,samples=70]({4*\x},{exp(\x-1)/(1-exp(-1))});
 \draw[curve,dashed]plot[domain=0:1,samples=70]({4*\x},{4*exp(4*(\x-1))/(1-exp(-4))});
 \node[below]at(4,0){$1$};\node at(1.3,1.6){高温};\node at(2.8,3.7){低温};
 \node at(2.4,-.85){$\delta Q=\dd E-\Omega\dd L_z$};
}
将已经求出的$\dd I/\dd T$代入图\ref{fig:rotcapacitya}对应的供热式，得到
\[
 C_\Omega=N\kb\left[\frac52-
 \frac{b_\Omega^2\ee^{-b_\Omega}}{(1-\ee^{-b_\Omega})^2}\right].
\]
弱旋转极限回到 $3N\kb/2$，强旋转时边壁附近的分布重排再贡献约 $N\kb$。该供热热容已经扣除了电机的机械交换。若装置改为与外界隔绝、角动量固定，升温时 $\Omega$ 也会改变，需要按新的约束重新求导。
\end{solution}

\begin{example}[重力气柱从低处泄流可以升温]
高 $H$、截面积 $S$ 的绝热刚性气柱中有 $N$ 个单原子理想气体粒子，温度 $T$ 均匀。高度 $h_0$ 处开小孔向真空缓慢泄流，气柱始终重新平衡，重力势能计入系统能量。写出 $T(N)$ 的微分关系；再求无限高气柱从底部泄流的温度粒子数关系。
\Diagram{fig:gravityleakq}{出口高度决定粒子携带的外场势能}{
 \draw(0,0)rectangle(2,3.2);\fill[white](1.94,1.05)rectangle(2.06,1.35);
 \draw[axis](2,1.2)--(4,1.2);\draw[guide](-.3,1.2)--(2,1.2);\node[left]at(-.3,1.2){$h_0$};
 \node at(3.2,1.9){$2\kb T+mgh_0$};\node at(1,-.5){势能零点在底部};
}
\end{example}
\begin{solution}
图\ref{fig:gravityleakq}每个逸出粒子平均携带 $\varepsilon_e=2\kb T+mgh_0$。余气单粒子平均总能量 $u(T)=3\kb T/2+mg\langle h\rangle$，固定 $H$ 的单粒子热容 $c(T)=\dd u/\dd T$。开放系统能量守恒为
\[
 u\dd N+Nc\dd T=\varepsilon_e\dd N,
 \qquad Nc\dd T=
 \left[\frac12\kb T+mg(h_0-\langle h\rangle)\right]\dd N.
\]
因此泄流是否冷却取决于出口能量相对于体内平均总能量的高低，不能只比较平动能。若 $h_0<\langle h\rangle-\kb T/(2mg)$，右侧方括号为负，因 $\dd N<0$，温度反而上升。
固定的重力场不显含时间，所以粒子穿孔时只把其所在高度的势能带出，不再额外计一次外场参数功。余气能量的全微分同时含 $u\dd N$ 和 $Nc\dd T$；前者是按体内平均能量减少粒子的贡献，后者才反映温度改变。底部逸出的粒子虽然平动能高于体内平动平均，却缺少体内平均具有的重力势能。比较总能量后，差值反而为负，这才使无限高气柱的底部泄流导致升温。

\Diagram{fig:gravityleaka}{无限高气柱底部逸出的总能量低于体内平均}{
 \node[draw,align=center]at(0,0){体内平均\\$5\kb T/2$};\node[draw,align=center]at(4.6,0){底部逸出\\$2\kb T$};\node at(2.3,-1){留下的粒子平均总能量升高};
}
无限高气柱有$mg\langle h\rangle=\kb T$、$c=5\kb/2$。从底部$h_0=0$泄流，图\ref{fig:gravityleaka}给
\[
 \frac52N\kb\dd T=-\frac12\kb T\dd N,
 \qquad T=T_0(N/N_0)^{-1/5}.
\]
分离变量后的初值积分为$5\int_{T_0}^{T}\dd\vartheta/\vartheta=-\int_{N_0}^{N}\dd N'/N'$，即$5\ln(T/T_0)=-\ln(N/N_0)$，给出上式。余气升温的原因是底部逸出粒子的总能量低于体内平均值。无限高度近似只适用于尺度高度远小于容器高度的阶段；温度升高使尺度高度增大后，应改用有限$H$的微分式描述。
\end{solution}

\begin{example}[光子工质的卡诺循环]
以平衡光子气体为工质，工作在 $T_h,T_c$ 之间。高温等温膨胀体积从 $V_A$ 到 $V_B$，随后可逆绝热到低温，再等温压缩和绝热返回。求其余两状态体积、两热源热量与循环功，验证卡诺效率。
\Diagram{fig:photoncarnotq}{光子气体的等温线是水平线}{
 \draw[axis](0,0)--(5.2,0)node[right]{$V$};\draw[axis](0,0)--(0,3.6)node[above]{$p$};
 \draw[curve](1,3)--(2,3);\draw[curve]plot[domain=2:3.90625](\x,{3*(2/\x)^(4/3)});
 \draw[curve](1.953125,1.2288)--(3.90625,1.2288);\draw[curve]plot[domain=1:1.953125](\x,{3/\x^(4/3)});
 \node[above]at(1,3){$A$};\node[above]at(2,3){$B$};\node[right]at(3.90625,1.2288){$C$};\node[left]at(1.953125,1.2288){$D$};\node at(3.9,3){$T_c/T_h=0.8$ 示意};
}
\end{example}
\begin{solution}
图\ref{fig:photoncarnotq}的绝热关系是 $VT^3=\text{常量}$，所以
\[
 V_C=V_B(T_h/T_c)^3,\qquad V_D=V_A(T_h/T_c)^3.
\]
高温等温段压强为 $aT_h^4/3$，内能也随体积增加。由 $Q=\Delta U+W$，
\[
 Q_h=\frac{4a}3T_h^4(V_B-V_A).
\]
低温段放热量大小 $Q_c=(4a/3)T_c^4(V_C-V_D)$，代入体积关系得到 $Q_c/Q_h=T_c/T_h$。
高温等温膨胀中，光子工质的熵增为$\Delta S_h=(4a/3)T_h^3(V_B-V_A)$。两条绝热段保持各端点的熵，因而低温压缩的熵减大小相同。低温体积差是高温体积差的$(T_h/T_c)^3$倍，而能量密度缩小为$(T_c/T_h)^4$，相乘后热量比只剩一次温度比。转到图\ref{fig:photoncarnota}的熵温平面，矩形面积为$(T_h-T_c)\Delta S_h$，等于两热源热量之差；原压强—体积图中的循环并非矩形。

\Diagram{fig:photoncarnota}{在熵温平面中不同工质具有相同的可逆循环结构}{
 \draw[axis](0,0)--(4.8,0)node[right]{$S$};\draw[axis](0,0)--(0,3.2)node[above]{$T$};\draw[curve](1,1)rectangle(3.6,2.5);\node[left]at(0,2.5){$T_h$};\node[left]at(0,1){$T_c$};\draw[axis](1.5,2.5)--(2.5,2.5);\draw[axis](3,1)--(2,1);\node at(2.3,1.8){循环功};
}
因此$W=Q_h(1-T_c/T_h)$，效率仍为卡诺值。图\ref{fig:photoncarnota}的矩形由两条等温线和两条等熵线组成；光子状态方程改变了各状态的体积及热量，不改变这一可逆循环的温度比效率。
\end{solution}

\section{磁性工质的熵与绝热降温}
\subsection{由两态分布接到广义功}
考虑$N$个相互独立的局域磁矩，磁矩大小为$\mu$，在外加磁感应强度$B>0$中有两个能级$\mp\mu B$。相互作用及内部磁场暂忽略，外场作为受控参数。令$x=\mu B/(\kb T)$，单个配分和为$2\cosh x$，于是
\[
 F_s=-N\kb T\ln(2\cosh x),\quad
 M=-\left(\frac{\partial F_s}{\partial B}\right)_T=N\mu\tanh x,
\]
\[
 U_s=-N\mu B\tanh x,\qquad
 S_s=N\kb[\ln(2\cosh x)-x\tanh x].
\]
这里$M$是样品沿场方向的总磁矩。所选系统能量包含磁矩与外加场的耦合能，磁场发生器及其场能在系统之外，相应热力学微分为$\dd U_s=T\dd S_s-M\dd B$。若改用单位体积磁化强度或把场源纳入系统，须同时调整能量与功的定义。

先有$\partial_Bx|_T=\mu/(\kb T)=x/B$，所以$M=-F_B=N\kb T\tanh x\,(x/B)=N\mu\tanh x$。从能级平均得到$U_s=-N\mu B\tanh x$后，用$S_s=(U_s-F_s)/T$得到上述熵。固定$B$时$\partial_Tx=-x/T$，因此$\partial_TU_s=-N\mu B\operatorname{sech}^2x(-x/T)$，再用$\mu B=\kb Tx$，得到
\[
 C_{B,s}=N\kb x^2\operatorname{sech}^2x.
\]
熵对$x$的导数为$N\kb[\tanh x-(\tanh x+x\operatorname{sech}^2x)]=-N\kb x\operatorname{sech}^2x$。再分别乘$x_T=-x/T$、$x_B=x/B$，得到
\[
 \dd S_s=\frac{C_{B,s}}T\dd T-\frac{C_{B,s}}B\dd B.
\]
对上述独立磁矩，可逆绝热过程满足$\dd T/T=\dd B/B$，即$B/T$保持不变，两个能级的占有比例也保持不变。推广到其他磁性物质或温度范围前，须重新计入相互作用及其他自由度的熵。
\ExpansionDiagram{fig:dep-th4-spins}{外场改变两能级间隔，热平衡占有比由$B/T$决定}{
\draw[->](0,0)--(0,3)node[above]{能量};\draw[thick](.6,.5)--(2.6,.5)node[right]{$-\mu B$};\draw[thick](.6,2.3)--(2.6,2.3)node[right]{$+\mu B$};\draw[<->](1.6,.6)--(1.6,2.2)node[midway,right]{$2\mu B$};\node[right]at(4,1.4){$N_+/N_-=\ee^{-2x}$};}
图\ref{fig:dep-th4-spins}中的能级标记以能量为准：$N_+$指上能级占有数，$N_-$指下能级。磁矩沿场方向对应低能级，不能把磁矩方向的正号与能级正号混用。

\subsection{晶格热容限制实际温降}
若样品还有与磁矩保持内部热平衡的晶格，晶格热容为$C_l(T)$且暂不依赖$B$，总熵微分为
\[
 \dd S=\frac{C_l+C_{B,s}}T\dd T-\frac{C_{B,s}}B\dd B.
\]
可逆绝热条件给
\[
 \left(\frac{\dd T}{\dd B}\right)_S
 =\frac TB\frac{C_{B,s}}{C_l+C_{B,s}}.
\]
晶格必须随磁矩一起降温，温降幅度因而小于独立自旋系统的$T\propto B$结果。若晶格与磁矩尚未来得及热平衡，则单一温度$T$已不足以描述过程，应另列两子系统的能量交换方程。

\begin{example}[有限晶格热容的绝热去磁]
样品的磁矩服从上述两态模型，晶格热容近似常数$C_l>0$。初态为$(B_i,T_i)$，其中$x_i=\mu B_i/(\kb T_i)$。样品与环境绝热并足够缓慢地减小外场到零，求理想模型给出的最终温度。特别取$x_i=1$、$C_l=N\kb$。
\ExpansionDiagram{fig:dep-th4-demag-q}{磁矩与晶格内部换热，外部绝热，仅改变受控磁场}{
\draw[dashed](-.4,-.4)rectangle(6,2);\draw(0,0)rectangle(2.2,1.5)node[midway]{磁矩};\draw(3,0)rectangle(5.4,1.5)node[midway]{晶格};\draw[<->](2.2,.7)--(3,.7);\node[above]at(2.7,2.1){总熵守恒};\draw[->](6.8,1.3)--(6.8,.1)node[below]{$B\downarrow$};}
\end{example}
\begin{solution}
图\ref{fig:dep-th4-demag-q}要求保存磁矩与晶格的总熵。晶格熵差是$C_l\ln(T_f/T_i)$。只要$T_f>0$，末场$B_f=0$时$x_f=0$，两个能级等概率，磁矩熵为$N\kb\ln2$。初末熵相等给
\[
 C_l\ln\frac{T_f}{T_i}+N\kb\ln2
 =N\kb[\ln(2\cosh x_i)-x_i\tanh x_i].
\]
所以
\[
 \frac{T_f}{T_i}=
 \exp\!\left\{\frac{N\kb}{C_l}
 [\ln\cosh x_i-x_i\tanh x_i]\right\}.
\]
括号在$x_i>0$时为负：其导数是$\tanh x_i-[\tanh x_i+x_i\operatorname{sech}^2x_i]=-x_i\operatorname{sech}^2x_i<0$，且原点为零。因此$0<T_f<T_i$，但有限正$C_l$使这个近似模型中的$T_f$不为零。
\ExpansionDiagram{fig:dep-th4-demag-a}{磁矩熵增加由晶格熵减少补偿；总熵保持不变}{
\draw[->](0,0)--(6,0)node[right]{过程};\draw[->](0,0)--(0,3.9)node[above]{$S$};\draw[thick](.5,1)--(5,2.2)node[right]{$S_s$};\draw[dashed](.5,2.4)--(5,1.2)node[right]{$S_l$};\draw[dotted](.5,3.4)--(5,3.4)node[right]{$S_{\rm total}$};}
图\ref{fig:dep-th4-demag-a}中磁矩熵的增加由晶格熵的减少补偿，线段仅表示变化趋势；温度由总熵方程求得。取$x_i=1,C_l=N\kb$，有$\ln\cosh1\simeq0.433781$、$\tanh1\simeq0.761594$，故$T_f/T_i\simeq\ee^{-0.327813}\simeq0.7205$。

若错误地仅保留磁矩子系统并令$T/B$不变，会预测$B\to0$时温度也趋零，遗漏了冷却晶格所需的熵转移。更低温时晶格热容并非常数，磁矩相互作用和内部场也可能不可忽略；届时应把实际$C_l(T)$积分写入总熵方程，而不是无限外推本题的常热容近似。
\end{solution}

\subsection{降温与系统能量降低不是同一个条件}
对只有磁矩的可逆绝热分支，$x$固定，$U_s=-N\mu B\tanh x$。减小$B$使$T$下降，却使负的$U_s$向零增加；由$\dd U_s=-M\dd B$，增加的能量来自受控外场参数的做功。温度反映占有分布相对能级间隔的尺度，不能脱离能量零点和外参变化，直接把降温解释为内能必定减少。

若要用这一工质进行制冷，应先等温磁化向环境排热，再绝热减场降温，随后与较冷负载接触吸热。各段固定的量不同，热量分别按$T\dd S$积分，磁场控制功按本节约定计入。只有把样品、负载与场源的边界交代清楚，才能比较一个完整制冷循环的能量与熵。

\section{习题与解答}
\begin{exercise}[等温混合大气的组分分层]
两种不发生反应的理想气体处于温度 $T$ 的重力平衡，摩尔质量为 $M_1,M_2$，地面分压为 $p_{10},p_{20}$。求各分压、摩尔分数与高度的关系，并讨论较重组分的比例。
\Diagram{fig:atmixq}{两种质量的粒子共享温度而尺度高度不同}{
 \draw(0,0)rectangle(2.6,3);\node at(1.3,.4){$p_{10},p_{20}$};\draw[guide](0,2)--(2.6,2);\node[right]at(2.8,2){$h$};\draw[axis](-.6,0)--(-.6,3.5);
}
\end{exercise}
\begin{solution}
各组分分别满足静力平衡，图\ref{fig:atmixq}中 $p_i(h)=p_{i0}\ee^{-M_igh/(RT)}$。因此 $x_i(h)=p_i(h)/[p_1(h)+p_2(h)]$，组分比为
\[
 \frac{x_2(h)}{x_1(h)}=\frac{p_{20}}{p_{10}}
 \exp\!\left[-\frac{(M_2-M_1)gh}{RT}\right].
\]
同一种外场中，不同组分的势能增量按各自摩尔质量计算；温度相同只意味着它们共享热能尺度 $RT$，并不使势能指数相同。总压 $p_1+p_2$ 通常不是用一个固定平均摩尔质量写成的单一指数。若 $M_2>M_1$，组分2的尺度高度更小，其相对比例随高度下降；若两摩尔质量恰相同，则指数相消，摩尔分数保持地面比例，恢复未发生质量分层的特例。

\Diagram{fig:atmixa}{较重组分的相对浓度随高度下降}{
 \draw[axis](0,0)--(4.7,0)node[right]{$h$};\draw[axis](0,0)--(0,2.8)node[above]{$x_2/x_1$};\draw[curve]plot[domain=0:4.3](\x,{2.3*exp(-.4*\x)});\node at(3,2.4){$M_2>M_1$};
}
图\ref{fig:atmixa}对应静态的分子扩散平衡，各组分按自己的摩尔质量分层。强湍流可以使较大高度范围内的组分比例近乎均匀，但这需要持续的混合作用，超出了此处静态模型的条件。
\end{solution}

\begin{exercise}[温度随高度线性下降的大气]
气体处于静力平衡，给定 $T(h)=T_0-\beta h>0$，$\beta>0$。摩尔质量 $M$、重力 $g$ 恒定。求 $p(h)$、$n(h)$。若局部满足理想气体的干绝热梯度，$\beta$ 应是多少？
\Diagram{fig:lapseq}{温度剖面由外部条件给定}{
 \draw[axis](0,0)--(4.6,0)node[right]{$T$};\draw[axis](0,0)--(0,3.2)node[above]{$h$};\draw[curve](4,0)--(1,2.7);\node[below]at(4,0){$T_0$};\node at(3,2.5){$T_0-\beta h$};
}
\end{exercise}
\begin{solution}
图\ref{fig:lapseq}不是全局等温平衡，因此不能直接写 $n\propto\ee^{-Mgh/(RT(h))}$。应从 $\dd p/p=-Mg\dd h/[R(T_0-\beta h)]$ 积分，得到
\[
 p=p_0\left(1-\frac{\beta h}{T_0}\right)^{Mg/(R\beta)},\qquad
 n=n_0\left(1-\frac{\beta h}{T_0}\right)^{Mg/(R\beta)-1}.
\]
积分中令$z=T_0-\beta h$，有$\dd h=-\dd z/\beta$，下限$h=0$变为$T_0$，上限$h$变为$T_0-\beta h$。所以
\[
 \ln\frac p{p_0}=\frac{Mg}{R\beta}
 \int_{T_0}^{T_0-\beta h}\frac{\dd z}{z}
 =\frac{Mg}{R\beta}[\ln z]_{T_0}^{T_0-\beta h}.
\]
数密度还要除以局部温度$T(h)$，因此幂指数比压强少1。若使$\beta\to0$，利用$\ln(1-\beta h/T_0)\simeq-\beta h/T_0$，所得压强恢复等温指数律；不能直接把变温度塞进等温指数的分母。

\Diagram{fig:lapsea}{微小气团绝热位移时焓变化与重力势能平衡}{
 \node[draw]at(0,0){$C_{p,m}\dd T$};\node at(2.3,0){$+$};\node[draw]at(4.2,0){$Mg\dd h$};\node at(2.1,-.8){$\dd H_m+Mg\dd h=0$};
}
按图\ref{fig:lapsea}，气团绝热位移时满足$\dd s=C_{p,m}\dd T/T-R\dd p/p=0$，所以$\dd p/p=(C_{p,m}/R)\dd T/T$。与静力式$\dd p/p=-Mg\dd h/(RT)$联立，约去$RT$后有
\[
 C_{p,m}\dd T=-Mg\dd h,\qquad
 \frac{\dd T}{\dd h}=-\frac{Mg}{C_{p,m}},\qquad
 \beta=\frac{Mg}{C_{p,m}}.
\]
这也等价于气团的摩尔焓增加与重力势能增加相抵。常热容时梯度为常数；若$C_{p,m}$随温度变化，应积分$\int C_{p,m}(T)\dd T=-Mg\int\dd h$。这种绝热分层或气团位移条件与全局等温平衡不同。
\end{solution}

\begin{exercise}[旋转容器内的气体分布]
圆柱容器半径 $R_c$、高 $L$，装有 $N$ 个质量 $m$ 的理想气体粒子，以角速度 $\Omega$ 匀速转动。温度 $T$ 均匀，忽略重力，气体相对容器达到平衡。求径向数密度和轴心、壁面压强之比。
\Diagram{fig:rotateq}{圆柱横截面中的半径与旋转方向}{
 \draw[curve](0,0)circle(1.8);\fill(0,0)circle(1.2pt);\draw[axis](0,0)--(1.5,.6)node[midway,above]{$r$};
 \draw[->](70:2.05)arc(70:160:2.05);\node at(-1.6,2.2){$\Omega$};\node at(2.8,0){$R_c$};
}
\end{exercise}
\begin{solution}
图\ref{fig:rotateq}在旋转系中有效势为 $E_p=-m\Omega^2r^2/2$，故 $n(r)=n_c\ee^{ar^2}$，$a=m\Omega^2/(2\kb T)$。归一化要用柱坐标体积元 $2\pi rL\dd r$，
\[
 N=2\pi Ln_c\int_0^{R_c}r\ee^{ar^2}\dd r
 =\frac{\pi Ln_c}{a}(\ee^{aR_c^2}-1).
\]
径向积分可令 $u=r^2$，使 $2r\dd r=\dd u$，原函数变为 $\ee^{ar^2}/(2a)$。若漏掉柱坐标的 $r$ 因子，就无法得到正确的归一化系数。气体向外受离心有效力，在旋转系的平衡式中由向内的压强合力抵消，故压强必须向外增大。轴心密度随总粒子数约束一起变化，升高角速度会使粒子向外缘集中，并不意味着可在保持原轴心密度的同时仍保持总量不变。

\Diagram{fig:rotatea}{密度向外增大，压强梯度提供向心加速度}{
 \draw[axis](0,0)--(4.7,0)node[right]{$r$};\draw[axis](0,0)--(0,2.8)node[above]{$n$};\draw[curve]plot[domain=0:4](\x,{.4*exp(.1*\x*\x)});\node at(2.7,2.5){$\dd p/\dd r=nm\Omega^2r$};
}
所以 $n_c=Na/[\pi L(\ee^{aR_c^2}-1)]$，而 $p(R_c)/p(0)=\ee^{aR_c^2}$。图\ref{fig:rotatea}的压强向外增大，其合力指向内侧，提供实验室系的向心加速度。$\Omega\to0$ 时用 $\ee^{aR_c^2}-1\simeq aR_c^2$，恢复均匀密度 $N/(\pi R_c^2L)$。
\end{solution}

\begin{exercise}[谐势阱的分布与热容]
$N$ 个经典独立粒子处于三维谐势 $E_p=(\kappa_xx^2+\kappa_yy^2+\kappa_zz^2)/2$ 中，温度为 $T$，空间无限，$\kappa_i>0$。求位置概率密度、各坐标均方值、总能量和热容。
\Diagram{fig:harmonicq}{各方向均有二次型束缚}{
 \draw[axis](-2.3,0)--(2.3,0)node[right]{$x$};\draw[axis](0,0)--(0,2.8)node[above]{$E_p$};\draw[curve]plot[domain=-1.8:1.8](\x,{.6*\x*\x});\node at(2,2.6){$\kappa_xx^2/2$};
}
\end{exercise}
\begin{solution}
图\ref{fig:harmonicq}每一方向的势能为正二次型，位置积分分别归一化，
\[
 \rho_P(x,y,z)=\frac{\sqrt{\kappa_x\kappa_y\kappa_z}}{(2\pi\kb T)^{3/2}}
 \exp\!\left[-\frac{\kappa_xx^2+\kappa_yy^2+\kappa_zz^2}{2\kb T}\right].
\]
每个坐标方向令$\alpha_i=\kappa_i/(2\kb T)$，已有高斯积分给
\[
 Z_i=\int_{-\infty}^{\infty}\ee^{-\alpha_ix_i^2}\dd x_i
 =\sqrt\pi\alpha_i^{-1/2},
\]
三个$Z_i$相乘就是上式的归一化分母。对参数求导，
\[
 -Z_i'=\int x_i^2\ee^{-\alpha_ix_i^2}\dd x_i
 =(\sqrt\pi/2)\alpha_i^{-3/2},
\]
因此
\[
 \langle x_i^2\rangle=-\frac{Z_i'}{Z_i}=\frac1{2\alpha_i}
 =\frac{\kb T}{\kappa_i}.
\]
每个势能项的平均值为$\kappa_i\langle x_i^2\rangle/2=\kb T/2$，总势能为$3N\kb T/2$。增大刚度缩小了空间宽度，经典平衡时每个二次项的平均能量仍只由温度决定。

\Diagram{fig:harmonica}{三个动能项和三个势能项各贡献半个热能尺度}{
 \node[draw,align=center]at(0,0){动能3项\\$3\kb T/2$};\node at(2.5,0){$+$};\node[draw,align=center]at(4.8,0){势能3项\\$3\kb T/2$};\node at(2.4,-.9){单粒子总能量 $3\kb T$};
}
图\ref{fig:harmonica}再加三个动能二次项，得到总能量 $U_{\rm tot}=3N\kb T$，热容 $3N\kb$。若只束缚 $d$ 个位置方向，其他方向由固定硬壁限制，则势能贡献为 $d\kb T/2$，总热容 $(3+d)N\kb/2$。量子低温下经典等分不再适用。
\end{solution}

\begin{exercise}[有限气柱高度的涨落]
单粒子在高 $H$ 的等温重力气柱中，令 $a_g=mg/(\kb T)$、$x=a_gH$。求平均高度与高度方差，证明 $\dd\langle mgh\rangle/\dd T=(mg)^2\operatorname{Var}(h)/(\kb T^2)$，并联系气柱的额外热容。
\Diagram{fig:heightvarq}{归一化的截断指数高度分布}{
 \draw[axis](0,0)--(4.8,0)node[right]{$h$};\draw[axis](0,0)--(0,3)node[above]{$w(h)$};\draw[curve]plot[domain=0:4](\x,{2.5*exp(-.5*\x)});\draw[guide](4,0)--(4,2.6);\node[below]at(4,0){$H$};
}
\end{exercise}
\begin{solution}
图\ref{fig:heightvarq}的归一化积分 $Z_h=\int_0^H\ee^{-a_gh}\dd h=(1-\ee^{-a_gH})/a_g$。对参数求导，$\langle h\rangle=-\partial\ln Z_h/\partial a_g$，再求一次导数得到方差：
\[
 \langle h\rangle=\frac1{a_g}-\frac H{\ee^x-1},\quad
 \operatorname{Var}(h)=\frac1{a_g^2}-\frac{H^2\ee^x}{(\ee^x-1)^2}.
\]
对显式的$\ln Z_h=\ln(1-\ee^{-a_gH})-\ln a_g$求导，得$\partial_{a_g}\ln Z_h=H/(\ee^{a_gH}-1)-1/a_g$。再求导时，$\partial_{a_g}[H/(\ee^{a_gH}-1)]=-H^2\ee^{a_gH}/(\ee^{a_gH}-1)^2$，$\partial_{a_g}(-1/a_g)=1/a_g^2$，因此$\partial_{a_g}^2\ln Z_h$等于上式方差。其概率意义已由位置配分积分的商法则证明：$-Z_h''/Z_h+(Z_h'/Z_h)^2=-\operatorname{Var}(h)$。这里高度积分的边界必须保持固定。

\Diagram{fig:heightvara}{参数导数把平均位置与位置涨落相联系}{
 \node[draw]at(0,0){$-\partial_{a_g}\ln Z_h$};\draw[axis](1.6,0)--(2.8,0);\node[draw]at(4.6,0){$\partial_{a_g}^2\ln Z_h$};\node at(2.3,-.9){平均高度的响应由方差决定};
}
由图\ref{fig:heightvara}，$\partial\langle h\rangle/\partial a_g=-\operatorname{Var}(h)$，且 $\dd a_g/\dd T=-mg/(\kb T^2)$，链式相乘即得题给关系。乘 $N$ 后，额外热容为 $N(mg)^2\operatorname{Var}(h)/(\kb T^2)\ge0$，与式\eqref{eq:columncapacity}一致。热容非负在此也由概率方差直接保证。
\end{solution}

\begin{exercise}[有限气柱的绝热外场变化（选做）]
单原子气体位于高 $H$ 的刚性气柱中，粒子数固定，初温 $T$。外场加速度 $g$ 缓慢改变，过程可逆绝热。用单粒子高度分布的配分积分推导 $T$ 对 $g$ 的局部变化率，并讨论无限高气柱极限。
\Diagram{fig:changegq}{外场改变会做功，即使容器不变形}{
 \draw(0,0)rectangle(2,3);\node at(1,1.5){$N,H$};\draw[axis](3,2.5)--(3,.7)node[midway,right]{$g\to g+\dd g$};\node at(1,-.5){$Q=0$};
}
\end{exercise}
\begin{solution}
图\ref{fig:changegq}中参数$g$改变，外界对粒子做广义功$Nm\langle h\rangle\dd g$，不能再写总能量不变。单粒子配分积分中动量项$\propto T^{3/2}$，位置项为$Z_h$。由式\eqref{eq:canonicalentropy}，再用$\langle E\rangle/T=3\kb/2+mg\langle h\rangle/T$，忽略与温度和$g$无关的常数，单粒子熵为
\[
 s=\kb\left[\frac32\ln T+\ln Z_h+\frac32+a_g\langle h\rangle\right].
\]
取 $T$ 不变求 $g$ 导数，利用 $\partial_{a_g}\ln Z_h=-\langle h\rangle$ 与 $\partial_{a_g}\langle h\rangle=-\operatorname{Var}(h)$，得到
\[
 \left(\frac{\partial S}{\partial g}\right)_T
 =-\frac{Nm^2g}{\kb T^2}\operatorname{Var}(h).
\]
具体地，固定$T$时$a_{g,g}=m/(\kb T)$，对位置熵两项求导有$\partial_g[\ln Z_h+a_g\langle h\rangle]=[-\langle h\rangle+\langle h\rangle-a_g\operatorname{Var}(h)]m/(\kb T)$。前两项抵消，乘$N\kb$即得上述结果。等熵微分还使用$S_T|_g=C/T$，这里$C$包含外场势能响应；只用平动热容会漏掉位置分布的储能变化。

\Diagram{fig:changega}{等熵约束联系温度响应与外场响应}{
 \node[draw]at(0,0){$(C/T)\dd T$};\node at(2.4,0){$+$};\node[draw]at(4.7,0){$S_g|_T\dd g$};\node at(2.4,-.8){总和为零};
}
按图\ref{fig:changega}，
\[
 \left(\frac{\dd T}{\dd g}\right)_S
 =\frac{Nm^2g\operatorname{Var}(h)}{\kb TC}.
\]
无限高气柱中方差$1/a_g^2$、$C=5N\kb/2$，故$\dd T/T=(2/5)\dd g/g$。从初值$(g_i,T_i)$积分，$[\ln\vartheta]_{T_i}^{T}=(2/5)[\ln g']_{g_i}^{g}$，即$T=T_i(g/g_i)^{2/5}$。外场增强使气柱压低，同时因广义做功而升温。
\end{solution}

\begin{exercise}[球面液滴与蒸气压的曲率修正（选做）]
不可压缩液体摩尔体积为 $v_l$，温度 $T$、表面张力 $\sigma$ 恒定。平面液面饱和蒸气压为 $p_\infty$。半径 $r$ 的液滴与理想蒸气平衡，忽略蒸气压改变量对液体压强的小贡献。推导液滴周围的平衡蒸气压。
\Diagram{fig:kelvinq}{球形界面的内外压强差}{
 \draw[curve](0,0)circle(1.5);\draw[axis](0,0)--(1.5,0)node[midway,above]{$r$};\node at(0,-.7){$p_l$};\node at(2.7,1){$p_v$};\node at(2.8,0){$p_l-p_v=2\sigma/r$};
}
\end{exercise}
\begin{solution}
图\ref{fig:kelvinq}的球面机械平衡给 $p_l-p_v=2\sigma/r$。相对于平面共存，液体化学势增量近似 $\Delta\mu_l=v_l\Delta p_l\simeq2\sigma v_l/r$。理想蒸气定温化学势变化为 $\Delta\mu_v=RT\ln(p_v/p_\infty)$。
液体侧定温积分为$\Delta\mu_l=\int_{p_\infty}^{p_l}v_l\dd p'=v_l(p_l-p_\infty)$，再忽略$p_v-p_\infty$，保留曲率压差。蒸气侧为$\Delta\mu_v=\int_{p_\infty}^{p_v}RT\dd p'/p'=RT[\ln p']_{p_\infty}^{p_v}$。两相都从平面共存参考态出发，化学势初值相等，故末态相平衡要求两增量相等。指数修正仅在不可压缩液体、理想蒸气及上述压差近似同时成立时适用。

\Diagram{fig:kelvina}{相平衡要求两相化学势增量相等}{
 \node[draw]at(0,0){$2\sigma v_l/r$};\node at(2.4,0){$=$};\node[draw]at(5,0){$RT\ln(p_v/p_\infty)$};
}
按图\ref{fig:kelvina}令两者相等，得
\[
 p_v=p_\infty\exp\!\left(\frac{2\sigma v_l}{rRT}\right).
\]
表面曲率提高了液体化学势，因而小液滴需要更高的蒸气压才能维持相平衡。液滴只有一个液气界面，压差为 $2\sigma/r$；肥皂泡的两层界面才给出 $4\sigma/r$。半径接近分子尺度时，连续介质表面张力模型失效。
\end{solution}

\begin{exercise}[共存线上潜热的局部变化]
不忽略两相摩尔体积，已知各相 $C_{p,m}$ 与 $v_T|_p$。从摩尔焓差 $L_m=h_g-h_l$ 出发，求沿共存线的 $\dd L_m/\dd T$；说明它何时化为前面的热容差近似。
\Diagram{fig:latentderivq}{沿共存曲线同时改变温度和压强}{
 \draw[axis](0,0)--(4.7,0)node[right]{$T$};\draw[axis](0,0)--(0,3)node[above]{$p$};\draw[curve](1,.4)to[out=30,in=240](3.7,2.5);\draw[axis](2,1)--(2.6,1.6);\node[right]at(3,1){$\dd p=p_s'(T)\dd T$};
}
\end{exercise}
\begin{solution}
图\ref{fig:latentderivq}的方向不是定压方向。对每一相，摩尔焓微分为 $\dd h=C_{p,m}\dd T+[v-Tv_T|_p]\dd p$，两相相减后除以 $\dd T$：
\[
 \frac{\dd L_m}{\dd T}=\Delta C_{p,m}
 +\left[\Delta v-T\Delta(v_T|_p)\right]\frac{\dd p_s}{\dd T}.
\]
沿共存线求导时，每个相的 $\dd p$ 都由同一个 $p_s'(T)\dd T$ 给出，所以压强项在两相相减后留下的是 $[v_g-v_l-T((v_g)_T-(v_l)_T)]p_s'$。其中的温度偏导按定压取，而整个 $L_m$ 的导数按沿共存线取，这两种导数有不同的约束。把它们都写成未注明条件的“对温度求导”，容易把第二项误删。最后再用 Clapeyron 式替换 $p_s'$，即可把修正完全用共存两相的性质表示。

\Diagram{fig:latentderiva}{潜热变化包含定压项与沿共存线的压强修正}{
 \node[draw]at(0,0){$\Delta C_{p,m}$};\node at(2.2,0){$+$};\node[draw]at(4.8,0){$[\Delta v-T\Delta v_T]p_s'$};
}
图\ref{fig:latentderiva}中第二项是沿共存线变压带来的修正。气相理想时 $v_g-T(v_g)_T|_p=0$，再忽略液体体积及热膨胀项，便得到 $\dd L_m/\dd T\simeq C_{p,m,g}-C_{p,m,l}=c_{v,g}+R-c_l$。前面的线性汽化热公式正是在这些近似下成立。
\end{solution}
